\documentclass[11pt,a4paper]{article}
\usepackage[T1]{fontenc}
\usepackage[utf8]{inputenc}
\usepackage{lmodern}
\usepackage[a4paper,margin=25mm,headheight=15pt]{geometry}
\usepackage{amsmath,amssymb,amsthm,mathtools}
\usepackage{microtype}
\usepackage{xcolor}
\usepackage{booktabs,array,longtable,enumitem}
\usepackage{tikz}
\usetikzlibrary{arrows.meta,positioning,calc,shapes.geometric}
\usepackage{fancyhdr}
\usepackage{hyperref}
\definecolor{ink}{HTML}{143C53}
\definecolor{accent}{HTML}{267A88}
\hypersetup{colorlinks=true,linkcolor=ink,urlcolor=accent,pdftitle={Noise-Rate-Activated Sensing: Explanatory Report},pdfauthor={Ruge Lin},pdfsubject={Global frequency estimation with a shared unknown noise direction},pdfkeywords={},bookmarksnumbered=true}
\pagestyle{fancy}
\fancyhf{}
\fancyhead[L]{\small\color{ink}Noise-Rate-Activated Sensing}
\fancyhead[R]{\small Explanatory report}
\fancyfoot[C]{\small\thepage}
\renewcommand{\headrulewidth}{0.3pt}
\setlength{\parindent}{0pt}
\setlength{\parskip}{0.55em plus 0.1em minus 0.1em}
\setlength{\emergencystretch}{2em}
\setlist{topsep=0.3em,itemsep=0.2em,parsep=0pt,leftmargin=1.6em}
\numberwithin{equation}{section}
\newtheorem{theorem}{Theorem}[section]
\newtheorem{proposition}[theorem]{Proposition}
\newtheorem{lemma}[theorem]{Lemma}
\theoremstyle{remark}
\newtheorem{remark}[theorem]{Remark}
\newcommand{\dd}{\mathrm d}
\DeclareMathOperator{\Tr}{Tr}
\DeclareMathOperator{\Var}{Var}
\DeclareMathOperator{\Cov}{Cov}
\newcommand{\E}{\mathbb E}
\newcommand{\R}{\mathbb R}
\allowdisplaybreaks[1]
\begin{document}
\begin{center}
  {\color{ink}\LARGE\bfseries Noise-Rate-Activated Sensing\par}
  \vspace{0.45em}
  {\large How changing noise strength changes the precision of an unknown frequency\par}
  \vspace{0.65em}
  {\normalsize Ruge Lin\par}
  \vspace{0.3em}
  {\small\parbox{0.94\textwidth}{\centering Thrust of Artificial Intelligence, Information Hub, The Hong Kong University of Science and Technology (Guangzhou), Guangzhou 511453, China}\par}
  \vspace{0.45em}
  {\small Explanatory report}
\end{center}
\vspace{0.4em}

How accurately can we read a clock when both its speed and the action of its
errors are unknown? In this report the clock is a quantum sensor. Its
\emph{phase} plays the role of a hand moving around a circle, and the frequency
is the speed we want to measure. Noise interrupts that motion. We know the
family of possible errors, but not which member of that family the device
has. The same unknown error acts throughout the experiment.

The main result compares two versions of one specified device. The first
has one fixed positive coupling to its noisy environment. The second may
choose between two fixed positive couplings to that same environment.
Neither setting turns the noise off. With only one coupling, the smallest
achievable worst-case mean-square error has order $1/n$, where $n$ is the
number of interaction cells. With the two-setting menu it has order $1/n^2$.
The clock interval, physical compensation and available boundary operations
are otherwise the same. The frequency is initially unknown on a fixed
nonzero interval, and the claim holds for every allowed error direction.

There are two different obligations in such a statement. To show that a
precision is \emph{attainable}, we must give a complete estimation procedure.
To show that it is \emph{optimal in order}, we must also rule out a better
order for every permitted procedure, including clever feedback schemes.
The report does both. Calibration, repeated trials, ambiguous phase readings
and unsuccessful outcomes all count toward the same physical duration.

\paragraph{The mechanism in ordinary language.}
A protected reference lets us detect when an error flips the sensor without
reading the phase we are trying to preserve. Two flips about the same unknown
axis cancel each other. The signal accumulates with opposite signs between
successive flips. Choosing the stronger coupling while the sensor is in the
flipped sector makes its return more likely; choosing the weaker one in the
unflipped sector helps retain a useful net signal. The complete error record
then tells us how each possible frequency would have affected the surviving
phase. Several progressively longer experiments locate the correct frequency.
This is a description of the construction; a separate bound is needed to show
why a single fixed coupling cannot achieve the same global precision order.

\paragraph{How to read the report.}
Section~\ref{intro:task} introduces the language through small examples.
Sections~\ref{micro:device} and \ref{phase:section} establish the main
finite-strength result. Sections~\ref{markov:continuous} and
\ref{markov:global} explain a separately proved continuous-time counterpart.
Sections~\ref{time:access} and \ref{time:timing} ask what changes if a pulse
can be applied during an interaction, and distinguish a known timing error
from hidden random jitter. The final sections state the limitations and
reassemble the physical conclusion. The prose explains the purpose of each
calculation; the equations specify exactly what is being claimed and why it
follows. No prior quantum mechanics or statistical estimation vocabulary is
assumed. Basic algebra is useful; slopes, averages and accumulated changes
are explained where their mathematical notation first appears.

\section{The estimation problem}
\label{intro:task}

\subsection{A two-level system can carry a phase}

A qubit is a system with two distinguished basis states, written
$|0\rangle$ and $|1\rangle$. The vertical bracket notation, called a ket,
denotes a column vector:
\begin{equation}
 |0\rangle=\begin{pmatrix}1\\0\end{pmatrix},\qquad
 |1\rangle=\begin{pmatrix}0\\1\end{pmatrix},\qquad
 |\psi\rangle=c_0|0\rangle+c_1|1\rangle,\qquad
 |c_0|^2+|c_1|^2=1.
 \label{intro:state}
\end{equation}
The coefficients $c_0,c_1$ are complex numbers called amplitudes.
A measurement in this basis returns 0 with probability $|c_0|^2$ and 1
with probability $|c_1|^2$. The last equation says that these probabilities
sum to one; this is what \emph{normalized} means. The basis states are
orthogonal: their inner product is zero, so this measurement can distinguish
them perfectly.

A complex amplitude has a magnitude and an angle. We write an angle factor
as $e^{i\vartheta}=\cos\vartheta+i\sin\vartheta$, where $i^2=-1$.
Multiplying the whole state by the same angle factor changes no measurement
probability. Changing the angle of one amplitude relative to the other,
however, changes what other measurements can reveal. That relative angle
is the phase used for sensing. For example,
\begin{equation}
 |\psi_\vartheta\rangle
 =\frac{|0\rangle+e^{i\vartheta}|1\rangle}{\sqrt2}
 \label{intro:phasestate}
\end{equation}
always gives equal probabilities in the 0/1 measurement, whatever
$\vartheta$ is. A suitable different measurement can distinguish its phases.

Operations on these vectors are represented by matrices. The three Pauli
matrices used throughout the report are
\begin{equation}
 X=\begin{pmatrix}0&1\\1&0\end{pmatrix},\qquad
 Y=\begin{pmatrix}0&-i\\i&0\end{pmatrix},\qquad
 Z=\begin{pmatrix}1&0\\0&-1\end{pmatrix}.
 \label{intro:pauli}
\end{equation}
Thus $X$ swaps the two basis states, $Z$ leaves $|0\rangle$ unchanged and
changes the sign of $|1\rangle$, and $Y$ swaps them with the phases $i$
and $-i$. The symbol $I$ is the identity matrix, which changes nothing.
A dagger means transpose followed by complex conjugation; conjugation
replaces $i$ by $-i$. A bra $\langle\psi|$ is the row vector
$|\psi\rangle^\dagger$. Hence $\langle\psi|\psi\rangle$ is the sum of the
squared amplitude magnitudes.

For measurements, $X,Y,Z$ also denote observables with outcomes $+1$ and
$-1$. An observable is a Hermitian matrix, meaning $O^\dagger=O$; its
eigenvalues are its possible outcomes. An eigenvector obeys
$O|v\rangle=\lambda|v\rangle$ and gives outcome $\lambda$ with certainty.
For the state in Eq.~\eqref{intro:phasestate},
\begin{equation}
 \langle X\rangle=\cos\vartheta,\qquad
 \langle Y\rangle=\sin\vartheta,\qquad
 \langle Z\rangle=0.
 \label{intro:quadratures}
\end{equation}
Angle brackets here denote the average measured outcome. Since the outcomes
are $\pm1$, a mean $m$ means probabilities $(1+m)/2$ and $(1-m)/2$.
The $X$ and $Y$ measurements are therefore two ways of reading the cosine
and sine of a phase. These are often called its two \emph{quadratures}.
They are performed on separately prepared systems; a single qubit does not
give both readouts for free.

\subsection{A density matrix keeps both phase and uncertainty}

The density matrix of a pure state is
$\rho=|\psi\rangle\langle\psi|$, a column times a row.
If different states are prepared with ordinary probabilities $p_j$, the
density matrix is their weighted average
$\rho=\sum_jp_j|\psi_j\rangle\langle\psi_j|$.
It has trace one, where $\Tr$ means the sum of the diagonal entries, and
is positive: $\langle v|\rho|v\rangle\ge0$ for every vector $|v\rangle$.
These properties ensure nonnegative measurement probabilities summing to one.

The difference between a quantum phase and an unrecorded random choice is
visible in a simple example:
\begin{equation}
 \rho_+=\frac12\begin{pmatrix}1&1\\1&1\end{pmatrix},
 \qquad
 \rho_{\mathrm{mix}}=\frac12\begin{pmatrix}1&0\\0&1\end{pmatrix}.
 \label{intro:mixture}
\end{equation}
The first is the state $|+\rangle=(|0\rangle+|1\rangle)/\sqrt2$.
The second is an equal random mixture of $|0\rangle$ and $|1\rangle$.
Both give 0 and 1 equally often in the basis measurement. An $X$ measurement
always gives $+1$ for $\rho_+$, but gives both signs equally often for
$\rho_{\mathrm{mix}}$. The off-diagonal entries preserve the phase
information that the mixture lacks.

In this notation, the mean outcome of observable $O$ is $\Tr(\rho O)$.
The probability of one outcome is $\Tr(\rho P)$, where $P$ is the projector
onto its outcome subspace. A projector keeps the component in that subspace
and removes the perpendicular component. These statements are the
\emph{Born rule}. They let us calculate the probabilities of every
experiment below, including experiments with partially lost phase information.

\subsection{Frequency turns the phase, and long turns can be ambiguous}

A Hamiltonian specifies how a quantum state changes in time. In units with
$\hbar=1$, evolution for time $t$ under a constant Hamiltonian $H$ is
$U(t)=e^{-iHt}$. The matrix exponential is defined by the same power series
as the ordinary exponential; it is a unitary matrix,
$U^\dagger U=I$, so evolution preserves normalization.
For $H=\omega Z/2$, its action on $|+\rangle$ is
\begin{equation}
 e^{-i\omega tZ/2}|+\rangle
 =\frac{e^{-i\omega t/2}|0\rangle+e^{i\omega t/2}|1\rangle}{\sqrt2}
 \ \sim\ \frac{|0\rangle+e^{i\omega t}|1\rangle}{\sqrt2}.
 \label{intro:clock}
\end{equation}
The symbol $\sim$ here means equality up to an unobservable overall phase.
The relative phase advances by $\omega t$. Frequency and coupling strength
have units of inverse time; phases are dimensionless and measured in radians.

A readout at known angle $\zeta$ uses
$O_\zeta=\cos\zeta\,X+\sin\zeta\,Y$. Its two outcome probabilities are
\begin{equation}
 \Pr\{y=\pm1\mid\omega,t,\zeta\}
 =\frac{1\pm\cos(\omega t-\zeta)}2.
 \label{intro:ramsey}
\end{equation}
This adjustable phase-sensitive readout is called an equatorial measurement.
An \emph{interrogation} means allowing the sensor to evolve and then
reading it. The displayed preparation, evolution and readout form a
Ramsey experiment; no further meaning of that name is needed here.

A long interrogation makes small frequency changes produce large phase
changes. But a full turn is $2\pi$: at a fixed duration $t$, frequencies
separated by $2\pi/t$ give exactly the same probabilities in
Eq.~\eqref{intro:ramsey}, even if the readout angle is varied.
The repeated oscillations of this probability are called fringes.
A long experiment can therefore resolve a small shift within one fringe
without identifying which fringe contains the true frequency.
Shorter measurements must first remove that ambiguity.

The device below includes a fixed compensating frequency $\omega_0$,
leaving the residual frequency $\epsilon=\omega-\omega_0$. This is the
known rotation already subtracted physically from the sensor's evolution.
Updating a classical estimate of $\epsilon$ changes what we know. Choosing
a different final readout angle changes what we measure. Neither action
changes the compensating Hamiltonian.

\subsection{Detect an error without reading the stored phase}

The sensor $S$ feels the signal and the noise. A protected ancillary qubit
$A$ acts as a reference and memory; it does not acquire an additional signal
phase. Environmental qubits $E$ interact with the sensor but are inaccessible.
The symbol $\otimes$ combines system spaces: $|0_S1_A\rangle$ means sensor
state 0 and ancilla state 1. Operators with different subscripts act on the
named systems. A partial trace removes one system from the mathematical
description while preserving every measurement probability on those remaining.
It does not make the removed system available to the controller.

The error family is
\begin{equation}
 N_\phi=\cos\phi\,X+\sin\phi\,Y,\qquad
 N_\phi|0\rangle=e^{i\phi}|1\rangle,\quad
 N_\phi|1\rangle=e^{-i\phi}|0\rangle.
 \label{intro:error}
\end{equation}
Every member swaps the two basis states, but leaves a different pair of
phases. The angle $\phi$ specifies which error occurs; it is fixed but
unknown throughout the experiment. We call it a transverse direction
because it reverses the signal generator:
\begin{equation}
 N_\phi^2=I,\qquad N_\phi ZN_\phi=-Z.
 \label{intro:erroralgebra}
\end{equation}
Two identical flips cancel. The reduced error channel is unchanged by
$\phi\mapsto\phi+\pi$, which changes only the sign of $N_\phi$; its
direction can therefore be taken modulo $\pi$.

Store a logical qubit in two physical qubits:
\begin{equation}
 |0_L\rangle=|0_S0_A\rangle,\quad
 |1_L\rangle=|1_S1_A\rangle,\quad
 |\psi_L\rangle=c_0|0_L\rangle+c_1|1_L\rangle.
 \label{intro:code}
\end{equation}
The subscript $L$ denotes this encoded two-dimensional space. Both states
have even parity, meaning that the two basis bits agree. The product
observable $Z_S\otimes Z_A$ then gives $+1$. A sensor error produces
\begin{equation}
 (N_\phi\otimes I_A)|\psi_L\rangle
 =c_0e^{i\phi}|1_S0_A\rangle+c_1e^{-i\phi}|0_S1_A\rangle,
 \label{intro:parity}
\end{equation}
which has odd parity: the bits disagree and the product observable gives
$-1$. Measuring parity distinguishes agreement from disagreement without
distinguishing the two amplitudes within either sector. It can reveal an
error without reading the signal phase. This error indication is called a
syndrome. It is obtained from the sensor and its reference, not from the bath.

A quantum \emph{channel} is the rule taking an input density matrix to the
accessible output. It can be written using Kraus operators as
\begin{equation}
 \rho\longmapsto\sum_rK_r\rho K_r^\dagger,\qquad
 \sum_rK_r^\dagger K_r=I.
 \label{intro:channel}
\end{equation}
The individual terms describe possible alternatives. Their traces are their
probabilities; the sum preserves total probability one. If a permitted
measurement also identifies the alternative $r$, we have an \emph{instrument}:
keep that label, and divide $K_r\rho K_r^\dagger$ by its trace to obtain
the conditional state. A Kraus label is not automatically observable.
Below, parity makes the relevant alternatives observable. An unrecorded
mixture and a collection of labelled conditional states are different
experimental resources.

\subsection{Error, probability and the resources we count}

A strategy specifies preparation, allowed operations, decisions based on
earlier observations, and the rule for turning the final record into an
estimate $\widehat\omega$. Adaptivity means applying one prescribed rule to
different observed histories. It does not give the controller the true
parameters. The complete record includes chosen settings, parity results,
terminal measurements and failure flags.

For a random number $W$, its expectation $\E W$ is its probability-weighted
average. Its variance is $\Var(W)=\E[(W-\E W)^2]$, the average squared
fluctuation around that mean. For independent observations with the same
variance, averaging $N$ of them divides that variance by $N$.
A conditional expectation $\E[W\mid R]$ uses only trials with the specified
record $R$, or equivalently holds that record fixed in the probability model.

The mean-square frequency error of a strategy $\mathcal P$ is
\begin{equation}
 \E_{\omega,\phi}^{\mathcal P}[(\widehat\omega-\omega)^2].
 \label{intro:mse}
\end{equation}
It averages \emph{all} its possible records, not only successful ones, and
has units of inverse time squared. To require a guarantee across an interval,
take the largest such error over its frequencies and all unknown directions.
Then choose the permitted strategy with the smallest possible guarantee.
This is the \emph{minimax risk}, written $R^*$. In the formal definition,
$\sup$ means the least upper bound (the largest value when one exists), and
$\inf$ means the greatest lower bound. Estimates may be biased: their
average need not equal the true frequency.

A fixed interval does not narrow as the experiment grows. A bound that is
\emph{uniform} over the interval and direction uses the same constant for
all those unknown values. The order statement $R^*(n)=\Theta(n^{-a})$
means that positive constants bound $n^aR^*(n)$ above and below for
sufficiently large $n$. $O(n^{-a})$ gives only an upper order bound, and
$o(n^{-a})$ denotes a term whose ratio to $n^{-a}$ tends to zero.
The hardware is held fixed in these limits.

For orientation, a pure $C/n$ mean-square law halves when the cell count
doubles, while $C/n^2$ quarters. The corresponding root-mean-square orders
are $n^{-1/2}$ and $n^{-1}$. Our theorem establishes these orders with
hardware-dependent constants, not these exact finite-budget ratios or a
particular numerical crossover between the two devices.

The cell duration is $h>0$ and the full physical duration is $T=nh$.
Ideal instruments have zero duration in this model, but their placement and
number are constrained. Initial preparation and readout exactly at the
deadline are free contacts. Earlier readouts, resets and interventions must
fit the contact budget. Calibration cells, failed trials and repetitions
all consume time. Storing an early estimate until the deadline creates no
additional sensing time. These are assumptions of the comparison, not a
claim that real laboratory control has no overhead.

\subsection{Local sensitivity and an unknown nuisance}

For a complete recorded outcome $x$, the likelihood $p(x\mid\omega,\phi)$
is the probability rule predicted at a candidate parameter pair.
A derivative is a local slope: $f'(u)$ tells how much $f$ changes per small
change in $u$. A partial derivative $\partial_\omega$ varies $\omega$
while holding the other parameter fixed. An integral later sums a continuously
changing contribution, such as elapsed signed time or an average over a
frequency interval.

The logarithmic derivatives
\begin{equation}
 \ell_\omega=\partial_\omega\log p,\qquad
 \ell_\phi=\partial_\phi\log p,\qquad
 J_{ij}=\E[\ell_i\ell_j]
 \label{intro:scores}
\end{equation}
are called scores and their Fisher information matrix. The logarithm is
natural throughout the report. A score measures the fractional change in an
outcome's probability. The matrix collects the size of those changes and
how similarly the two unknowns change the data.

The direction $\phi$ is a nuisance parameter: it affects the observations
but is not the quantity we wish to report. If a frequency change can be
imitated by changing $\phi$, its apparent sensitivity is not all useful.
The remaining local information is
\begin{equation}
 \mathcal I_\omega
 =\min_c\E[(\ell_\omega+c\ell_\phi)^2]
 =J_{\omega\omega}-\frac{J_{\omega\phi}^2}{J_{\phi\phi}}
 \quad\text{if }J_{\phi\phi}>0.
 \label{intro:efficient}
\end{equation}
The minimization asks how small the observable change can be made by
combining the frequency change with a direction change. It is a comparison
of possible unknown parameters, not a controller action. The coefficient
$c$ has units of time. The minimization form remains meaningful when the
division formula is unavailable. Trials sharing one direction must combine
their full matrices before eliminating that same unknown direction.

Quantum Fisher information bounds the information obtainable by any
measurement of a quantum state. This helps prove an impossibility result for
all readouts. Both forms of Fisher information describe \emph{local}
distinguishability near a parameter value. The repeated-fringe example
explains why large local information does not itself give a globally accurate
estimate. The lower bounds below turn uniform information bounds into error
bounds using a carefully chosen average over possible frequencies. The upper
bounds instead construct actual estimators and charge every part of their
operation.


\section{A finite interaction device and its precision limit}
\label{micro:device}

The central comparison can be made without taking a weak-coupling or
infinitely fast measurement limit. A sensor meets a sequence of fresh bath
qubits. Each encounter lasts the same positive time $h$, and the coupling
strength remains finite as the experiment grows. The controller may act
between encounters. Its additional option in the more powerful design is to
select one of two positive strengths for the next encounter.

The target frequency is $\omega$, measured in inverse time. A fixed nominal
compensation $\omega_0$ leaves the residual frequency
$\epsilon=\omega-\omega_0$. The unknown transverse direction is described by
an angle $\phi$, shared by every encounter, setting and repeated trial.
On a sensor qubit, let $X_S,Y_S,Z_S$ denote the Pauli matrices, and put
\begin{equation}
 N_\phi=\cos\phi\,X_S+\sin\phi\,Y_S.
 \label{micro:axis}
\end{equation}
This operator obeys $N_\phi^2=I$ and $N_\phi Z_SN_\phi=-Z_S$.
An error about this direction reverses the sign of the signal generator;
two such errors cancel. The difficulty is that the direction itself is
unknown and continuously valued. For the reduced experiment, $\phi$ and
$\phi+\pi$ describe the same noise axis.

\subsection{What is fixed, what is controlled, and what is counted}

During a cell with setting $j\in\{a,b\}$ and coupling $g_j>0$, the joint Hamiltonian is
\begin{equation}
 H_j(\epsilon,\phi)
 =\frac{\epsilon}{2}Z_S\otimes I_E+g_jN_\phi\otimes X_E.
 \label{micro:hamiltonian}
\end{equation}
The subscript $E$ denotes that cell's environmental qubit. Units are chosen
so that $\hbar=1$: $H_j$, $g_j$ and $\epsilon$ have units of inverse time,
whereas $g_jh$ is dimensionless. All bath qubits start in $|0\rangle$ at time
zero. Each is used once and retained afterward. The controller cannot
measure or manipulate them. Thus a finite experiment uses a finite bath,
whose size grows with the number of cells.

The controller has protected ancillary storage and ideal, instantaneous
instruments at the boundaries. An instrument is an operation that can both
change the quantum state and produce a classical outcome. Its subsequent
choices may depend on the outcomes already recorded. The control rule is
fixed independently of the unknown parameters: learning from data is
allowed, but the controller is not given the true values of $\omega$ or
$\phi$.

A budget of $n$ cells has hard duration $T=nh$. Initial preparation and
readout at $T$ are free. There are at most $n-1$ nonterminal boundary
contacts. Several consecutive operations at the same boundary count as
one combined instrument: a parity measurement, a completed trial's
readout, a fresh preparation and a choice of the next coupling can be
combined there. A readout before $T$ is a counted contact. Calibration,
repetitions, rejected records and unsuccessful trials all consume their
cells and contacts. An estimate obtained early can be held in protected
storage until the deadline, but the charged duration remains $nh$.

There are no controls inside a cell, changes of cell duration, resets of
the collision clock, or noise-free signal-bearing waiting periods. The
physical compensation $\omega_0$ stays fixed. A frequency estimate may
change a later measurement basis; it does not change the Hamiltonian in
Eq.~\eqref{micro:hamiltonian}.

Choose two specified couplings satisfying
\begin{equation}
 0<g_a<g_b,\qquad g_bh<\frac{\pi}{4}.
 \label{micro:couplings}
\end{equation}
The fixed-device benchmark uses one specified $g\in\{g_a,g_b\}$ in every
cell of every pooled trial. The two-setting device may choose $g_a$ or
$g_b$ before each cell through classical recorded control. A recorded
random choice between strengths is already access to the two-setting
menu. Both classes have the same boundary-control and storage resources.
Neither coupling decreases as $n$ grows, and neither is zero.

\subsection{The global statement}

The frequency is promised to lie in a fixed interval
$\Omega=[\omega_0-E,\omega_0+E]$, with $E>0$ independent of $n$.
Its sufficient size will be specified below. For a control protocol
$\mathcal P$ and estimator $\widehat\omega$, the risk is the mean squared
error, averaged over every recorded outcome. The optimal worst-case risk
at one coupling is
\begin{equation}
 R_g^*(n;\Omega)
 =\inf_{\mathcal P,\widehat\omega}
   \sup_{\omega\in\Omega,\,\phi}
   \mathbb E_{\omega,\phi}^{\mathcal P}
       [ (\widehat\omega-\omega)^2 ].
 \label{micro:risk}
\end{equation}
The same definition with the two-setting control class gives
$R_{\{g_a,g_b\}}^*(n;\Omega)$. The infimum includes biased estimators and
all permitted adaptive strategies, not just the constructions described
in this report. Squared frequency error has units of inverse time squared.
Clipping any estimate to $\Omega$ cannot increase its risk.

\begin{theorem}[Global finite-strength separation]
\label{micro:theorem}
For every fixed pair in Eq.~\eqref{micro:couplings}, there is a fixed,
nonzero bandwidth $E$ on which, uniformly over the shared unknown
transverse direction,
\begin{equation}
 \boxed{\quad
 R_g^*(n;\Omega)=\Theta(n^{-1})\quad(g=g_a\text{ or }g_b),
 \qquad
 R_{\{g_a,g_b\}}^*(n;\Omega)=\Theta(n^{-2}).\quad}
 \label{micro:separation}
\end{equation}
The clock, strengths and compensation remain fixed throughout this
comparison.
\end{theorem}
Here $\Theta(n^{-k})$ means that positive, finite constants bound the risk
above and below by multiples of $n^{-k}$ for all sufficiently large $n$.
The constants may depend on the fixed hardware and interval. As $h$ is
fixed, the corresponding laws are inverse physical time and inverse
physical time squared. This theorem establishes orders of risk; it does
not identify the optimal global constants or the largest possible band.

Three ingredients establish the comparison. An exact one-cell experiment
gives the inverse-$n$ upper bound and provides paid coarse localization.
A direction-changing comparison curve proves that every fixed-device
strategy faces the inverse-$n$ obstruction. Finally, the two-strength
construction retains records whose accumulated phase is uniformly
sensitive to frequency and resolves their ambiguity across several
scales. The latter construction is developed in the next part of the
report.

\subsection{The exact instrument of one cell}

We need both parts of a cell's output: what the parity measurement reports,
and what happens to the stored phase after that report. An average noise
channel alone would discard this distinction. The following calculation
provides the probabilities and the conditional operations together.

The two terms in Eq.~\eqref{micro:hamiltonian} anticommute: multiplying them
in opposite orders changes the sign, so their cross terms cancel when the
Hamiltonian is squared. Therefore
$H_j^2=R_j(\epsilon)^2I$, where
\begin{equation}
 R_j(\epsilon)=\sqrt{g_j^2+\epsilon^2/4},\qquad r_j=R_j(\epsilon)h.
\end{equation}
The symbol $R_j(\epsilon)$ here is a frequency scale, distinguished by
its argument and absence of a star from the risk $R_g^*(n;\Omega)$.
Exponentiating the Hamiltonian gives the exact unitary
\begin{equation}
 U_j=\cos r_j\,I
     -i\frac{\sin r_j}{R_j}
       \left(\frac{\epsilon}{2}Z_S\otimes I_E
                        +g_jN_\phi\otimes X_E\right).
 \label{micro:unitary}
\end{equation}
The operator norm, meaning the largest factor by which a matrix can stretch
a vector, is
$\|H_j\|=R_j(\epsilon)$. No short-time expansion has entered this formula.

Tracing out a bath input $|0\rangle$ produces a quantum channel
$\rho\mapsto\sum_{r=0}^1K_{r,j}\rho K_{r,j}^\dagger$ with operators
\begin{equation}
 K_{0,j}=a_jI-ib_jZ_S,\qquad K_{1,j}=-ic_jN_\phi,
 \qquad
 a_j=\cos r_j,\quad
 b_j=\frac{\epsilon}{2R_j}\sin r_j,\quad
 c_j=\frac{g_j}{R_j}\sin r_j.
 \label{micro:kraus}
\end{equation}
These real coefficient functions satisfy $a_j^2+b_j^2+c_j^2=1$. The term
``Kraus operators'' denotes this representation of an unobserved
system--environment interaction; it does not imply environmental access.

One protected ancillary qubit makes the alternatives observable without
reading out the stored phase. Encode
$|0_L\rangle=|00\rangle$ and $|1_L\rangle=|11\rangle$, with the first
qubit the sensor. The subscript $L$ denotes this two-dimensional logical
space. The first operator preserves the parity of the two qubits, while
the second reverses it. A parity measurement distinguishes these cases
without distinguishing $|0_L\rangle$ from $|1_L\rangle$. Because
\begin{equation}
 K_{1,j}^\dagger K_{1,j}=p_jI,\qquad K_{0,j}^\dagger K_{0,j}=q_jI,
 \qquad
 p_j(\epsilon)=\frac{g_j^2}{R_j^2}\sin^2(R_jh),\quad q_j=1-p_j,
 \label{micro:probabilities}
\end{equation}
the flip and no-flip probabilities are independent of the logical state.
They are also independent of $\phi$, although generally dependent on
$\epsilon$.

On the band used below, the normalized no-flip branch is a phase rotation:
\begin{equation}
 \frac{K_{0,j}}{\sqrt{q_j}}
  =\exp[-i\theta_j(\epsilon)Z_S/2],\qquad
 \theta_j(\epsilon)
  =2\arctan\left[\frac{\epsilon}{2R_j}\tan(R_jh)\right].
 \label{micro:theta}
\end{equation}
Here $\arctan$ is the inverse tangent taking values between $-\pi/2$ and
$\pi/2$. This is the continuous phase branch with $\theta_j(0)=0$.
The flip branch is $N_\phi$ up to an irrelevant scalar phase.
A phase is dimensionless; its derivative with respect to frequency has
units of time. In particular,
\begin{equation}
 d_j=\theta_j'(0)=\frac{\tan(g_jh)}{g_j},\qquad j\in\{a,b\}.
 \label{micro:derivative}
\end{equation}
A prime denotes differentiation with respect to residual frequency, not
time. Although $d_j>h$, this conditional derivative is not a gain obtained
without paying for failure: at compensation,
$q_jd_j^2=h^2\sin^2(g_jh)/(g_jh)^2\le h^2$.

\subsection{One fixed band for calibration and coherent accumulation}
\label{micro:band}

The frequency promise must do three jobs. It must keep the one-cell response
invertible, keep the flip probabilities positive and separated, and preserve the
useful slope of long recorded phases. We now impose sufficient margins for
all three jobs on the same interval. Their precise fractions are convenient
proof choices; the physical point is that one positive-width interval works
for every later budget. Choosing this interval does not require knowing
the actual frequency within it.

Write $\chi_j=g_jh$ and choose
$\chi_b<r_*<\pi/4$. Begin with
\begin{equation}
 0<E\le\frac{2}{h}\sqrt{r_*^2-\chi_b^2}.
 \label{micro:band1}
\end{equation}
Then $R_j(\epsilon)h\le r_*$ for both settings and every
$|\epsilon|\le E$. In particular the complete Hamiltonian is bounded by
$r_*/h$, not just its coupling term. A specified hardware strength limit
must cover this complete bound.

The setting-dependent functions are $p_j$, $q_j$ and $\theta_j$. Define the positive time scale
\begin{equation}
 \mu_0=
 \frac{p_b(0)q_a(0)d_a-p_a(0)q_b(0)d_b}{p_a(0)+p_b(0)}.
 \label{micro:mu}
\end{equation}
Its meaning will be the average sensitivity gained per cell when the
strength is chosen from the observed parity. A ``reward'' below is this
signed sensitivity score, not a payment or another physical resource.
Its positivity can already be
checked algebraically. In the trigonometric expressions below,
$\tan x=\sin x/\cos x$, $\cot x=\cos x/\sin x$, and $\csc x=1/\sin x$.
At compensation,
\begin{equation}
 \frac{q_j(0)d_j}{p_j(0)}=h\frac{\cot\chi_j}{\chi_j},
 \qquad
 \frac{\dd}{\dd\chi}\frac{\cot\chi}{\chi}
 =-\frac{\chi\csc^2\chi+\cot\chi}{\chi^2}<0.
 \label{micro:driftpositive}
\end{equation}
Since $\chi_a<\chi_b$, the numerator in Eq.~\eqref{micro:mu} is
strictly positive.

At nonzero residual frequency define the frozen-derivative mean
\begin{equation}
 \mu^0(\epsilon)=
 \frac{p_b(\epsilon)q_a(\epsilon)d_a
              -p_a(\epsilon)q_b(\epsilon)d_b}
      {p_a(\epsilon)+p_b(\epsilon)}.
\end{equation}
The superscript $0$ means that the rewards $d_a,d_b$ are evaluated at
zero, while the probabilities retain their actual frequency dependence.
Shrink $E$, if necessary, until throughout the band
\begin{align}
 &0<p_a(\epsilon)<p_b(\epsilon)<\tfrac12,
 \qquad \mu^0(\epsilon)\ge\tfrac34\mu_0,
 \label{micro:band2}\\
 &\max_{j\in\{a,b\}}|\theta_j'(\epsilon)-d_j|
                         \le\tfrac14\mu_0.
 \label{micro:band3}
\end{align}
These are conditions on known smooth functions of one real variable.
Their strict values at zero and continuity guarantee some positive
bandwidth. Fix that bandwidth once. It is not narrowed as the cell
budget grows. The conditions also explain why constants need not remain
uniform if the two couplings approach one another: their positive drift
then disappears.

\subsection{Calibration that uses every outcome}
\label{micro:calibration}

Fix one coupling and suppress its setting label, writing $g,R,r,a,b,c,p,q,\theta,H,U,K_0,K_1$
for the corresponding cell quantities. Prepare $|+_L\rangle=(|00\rangle+|11\rangle)/\sqrt2$ for one cell.
After measuring parity, measure the logical sine quadrature $Y_L$ if
parity is even. This is the Pauli $Y$ measurement in the logical basis:
its outcomes are $\pm1$, and their mean for a phase-rotated
$|+_L\rangle$ is the sine of that phase. Record its outcome
$W\in\{-1,+1\}$. If parity is odd,
record $W=0$. $W$ is a classical score, distinct from a Pauli
operator. Equation~\eqref{micro:kraus} gives
\begin{equation}
 \Pr\{W=0\}=p(\epsilon),\qquad
 \Pr\{W=\pm1\}=\frac12(a\pm b)^2
              =\frac{q(\epsilon)\pm f(\epsilon)}2,
 \label{micro:caliblikelihood}
\end{equation}
where
\begin{equation}
 f(\epsilon)=\mathbb E W
       =\frac{\epsilon\sin(2Rh)}{2R}.
 \label{micro:F}
\end{equation}
Every probability is independent of the unknown direction. The odd
outcome is part of the experiment, not an event removed from its cost or
likelihood. All three probabilities are strictly positive on the chosen
band, because $|b|<\sin r<\cos r=a$ and $p>0$.

For a concrete trial, take $gh=\pi/6$ and $\epsilon=0$. Then $p=1/4$,
$a=\sqrt3/2$ and $b=0$: the probabilities of $0,+1,-1$ are
$1/4,3/8,3/8$. The mean score is zero for every unknown direction.
A small positive residual frequency makes $+1$ more likely than $-1$;
a negative one reverses the imbalance. Repetition measures this imbalance.
The example illustrates the readout, not the two-strength advantage.

The response $f$ is monotone with a uniform inverse. Differentiation
shows
\begin{equation}
 f'(\epsilon)
 =h\left[
    \frac{g^2}{R^2}\frac{\sin(2r)}{2r}
   +\frac{\epsilon^2}{4R^2}\cos(2r)\right].
 \label{micro:Fprime}
\end{equation}
The two weights are nonnegative and add to one. For $0<x<\pi/2$,
$\sin x/x\ge\cos x$; indeed the derivative of
$\sin x-x\cos x$ is $x\sin x\ge0$, and its value at zero is zero.
It follows that
\begin{equation}
 \kappa\le f'(\epsilon)\le h,
 \qquad \kappa=h\cos(2r_*)>0.
 \label{micro:kappa}
\end{equation}
Thus distinct frequencies in the whole promised band give distinct mean
scores, with a slope that never becomes arbitrarily small there.

Repeat this complete preparation and readout for $N$ cells. Fresh bath
partners and re-preparation make the scores independent. With
$\overline W=N^{-1}\sum_{j=1}^NW_j$, define
\begin{equation}
 \widehat\epsilon_N=
 f^{-1}\!\left(\operatorname{clip}_{[f(-E),f(E)]}
                          (\overline W)\right).
 \label{micro:calibestimator}
\end{equation}
The physical-frequency estimate is $\widehat\omega=\omega_0+\widehat\epsilon_N$.
Here clipping projects a number to the nearest point in the stated
interval. It cannot increase its distance from the true mean. The
variance $\operatorname{Var}(W)$ is the average squared deviation from
that mean. Since $W\in[-1,1]$, it is at most one. Independence makes the
variance of the sample mean $\operatorname{Var}(W)/N$: cross terms from
different centered scores average to zero. The inverse slope bound therefore
yields
\begin{equation}
 |\widehat\omega-\omega|
 \le\kappa^{-1}|\overline W-f(\epsilon)|,
 \qquad
 \sup_{\omega\in\Omega,\phi}
 \mathbb E(\widehat\omega-\omega)^2
 \le\frac1{N\kappa^2}.
 \label{micro:calibrisk}
\end{equation}
The same bounded scores give the uniform concentration bound
\begin{equation}
 \Pr\{|\widehat\omega-\omega|\ge\delta\}
 \le2\exp[-N\kappa^2\delta^2/2],\qquad\delta>0.
 \label{micro:calibtail}
\end{equation}
This bounds the probability that the frequency error exceeds a chosen
tolerance $\delta$. The bound decreases exponentially with the number of
trials; it is stronger than a statement about average error alone.
For confidence failure probability at most $\eta\in(0,1)$, it is enough
to take
\begin{equation}
 N=\left\lceil\frac{2\log(2/\eta)}{\kappa^2\delta^2}\right\rceil.
 \label{micro:calibcost}
\end{equation}
The ceiling brackets mean rounding up to a whole number of cells.
All logarithms here and below are natural unless explicitly stated
otherwise. This is an ordinary concentration inequality applied to an
exact microscopic likelihood; it is not a weak-coupling approximation.

As a stand-alone experiment, these $N$ cells use $N-1$ nonterminal
contacts and terminal readout at $Nh$. As the prefix of a longer run,
they use $N$ counted contacts: the readout and next preparation at $Nh$
are then before the final deadline. The prefix also costs time $Nh$,
and its failure probability remains part of the later estimator's risk.
Taking $N=n$ proves the fixed-device upper bound. Using a shorter,
explicitly charged prefix supplies the coarse interval needed for the
two-setting estimator.

\subsection{Why arbitrary adaptation cannot beat the fixed-device order}
\label{micro:converse}

The obstruction is not that the preceding calibration happens to use
independent cells. It applies to every allowed boundary strategy,
including strategies that try to learn $\phi$ while preserving quantum
coherence in memory. To demonstrate this, consider a mathematical curve
along which a frequency change is accompanied by a suitable change in
the unknown direction. A worst-case guarantee over all $(\omega,\phi)$
must also hold on this curve. No physical relationship between the two
unknowns is imposed, and the curve is not supplied to the controller.

Fix one coupling $g$ and use the coefficient functions $a,b,c$ from
Eq.~\eqref{micro:kraus}. Put
\begin{equation}
 p=c^2,\qquad q=a^2+b^2,\qquad s=ab'-ba',\qquad
 \phi(\epsilon)=\phi_0+\int_0^\epsilon\frac{s(u)}{p(u)}\,\dd u.
 \label{micro:curve}
\end{equation}
The integral is well-defined because, on the entire band,
\begin{equation}
 p\ge p_{\min,g}\equiv(gh)^2
                 \left(\frac{\sin r_*}{r_*}\right)^2>0.
 \label{micro:pmin}
\end{equation}
This inequality follows from
$p=(gh)^2(\sin r/r)^2$ and the decrease of $\sin r/r$ on the
positive range in use. Angles in Eq.~\eqref{micro:curve} can be read
modulo $\pi$; the smooth real-valued expression is a convenient lift.

A dot now means the total derivative along this curve: both $\epsilon$
and $\phi(\epsilon)$ vary. Direct multiplication, using
$N_\phi\partial_\phi N_\phi=iZ_S$, gives
\begin{align}
 \sum_{r=0}^1K_r^\dagger\dot K_r
     &=i(-s+p\phi')Z_S=0,
 \label{micro:orthogonality}\\
 \sum_{r=0}^1\dot K_r^\dagger\dot K_r
     &=A(\epsilon)I,
 \qquad A=a'^2+b'^2+c'^2+\frac{s^2}{p}.
 \label{micro:A}
\end{align}
The first identity is the decisive cancellation. The direction changes
at precisely the rate needed to remove the coherent accumulation of the
channel derivative. Here is why this matters. If small changes produced by
successive cells all point in the same direction in state space, their
sizes can add before squaring, giving a term of order $n^2$. If the changes
are perpendicular, their squared sizes add instead, giving order $n$.
The zero-overlap identity enforces the latter behavior along the comparison
curve, even when the controller uses the earlier record.
Equivalently, $s=q\theta'/2$, and
\begin{equation}
 4A=\frac{p'^2}{pq}+\frac{q\theta'^2}{p}.
 \label{micro:Aequivalent}
\end{equation}
The first term reflects changes in the branch probabilities; the second
accounts for phase sensitivity and the direction variation. Dropping the
probability term would not be a valid off-compensation calculation.

There is a simple uniform bound. At fixed $\phi$,
$\partial_\epsilon H=Z_S\otimes I_E/2$. Differentiating the unitary
as an integral over the evolution gives
$\|\partial_\epsilon U\|\le h/2$. Applied to the bath input, this
implies $A_0=a'^2+b'^2+c'^2\le h^2/4$. The ordinary
Cauchy--Schwarz inequality, which bounds the size of an inner product
by the product of the two vector lengths, gives
$s^2\le q(a'^2+b'^2)\le qA_0$. Therefore
\begin{equation}
 4A\le\frac{h^2}{p}\le\frac{h^2}{p_{\min,g}}.
 \label{micro:Abound}
\end{equation}

To see why the bound survives feedback, represent one cell by the
isometry $V=\sum_rK_r\otimes|r\rangle$, where the extra orthogonal
labels retain a mathematical purification of the discarded bath.
An isometry preserves vector lengths, so $V^\dagger V=I$.
A purification enlarges the state space while reproducing the physical
state when the extra registers are ignored. It can only increase the
information available for an upper bound. Similarly purify boundary
instruments and keep all classical records in orthogonal registers.
The controller's instrument, including its conditional branches, has no
explicit parameter derivative.

If $|\Psi\rangle$ is the accumulated purified state, the derivative after
a cell is $\dot V|\Psi\rangle+V|\dot\Psi\rangle$. Expanding its squared
length gives
\begin{equation}
 \|\dot V\Psi+V\dot\Psi\|^2
 =\|\dot V\Psi\|^2+\|\dot\Psi\|^2
   +2\operatorname{Re}\langle\dot\Psi|V^\dagger\dot V|\Psi\rangle.
\end{equation}
Here $\operatorname{Re}$ takes the real part of a complex number.
The last, cross term vanishes because $V^\dagger\dot V=0$ from
Eq.~\eqref{micro:orthogonality}. Each cell consequently adds at most $A$
to the derivative norm squared. Induction gives an upper bound $nA$.
For a normalized pure state, the quantum Fisher information is
$4(\langle\dot\Psi|\dot\Psi\rangle-
 |\langle\Psi|\dot\Psi\rangle|^2)$, which is at most four times the
first term. Every actual final measurement has classical Fisher
information $J_\epsilon$ no larger than this. Hence
\begin{equation}
 J_\epsilon\le F_\epsilon^{\mathrm{out}}
          \le4nA(\epsilon)
          \le\frac{nh^2}{p_{\min,g}}.
 \label{micro:adaptivebound}
\end{equation}
The classical information is the mean squared likelihood score,
$J_\epsilon=\mathbb E[(\partial_\epsilon\log p(x|\epsilon))^2]$;
$F_\epsilon^{\mathrm{out}}$ denotes the corresponding maximal quantum
information before final measurement. The variable $x$ is the complete
outcome, including any flags, settings and stopping decisions.

This proof charges the whole adaptive experiment and keeps failures.
Resetting the sensor, learning the direction, branching on earlier data
or storing a result early does not remove the derivative cancellation.
A separate curve is allowed for each fixed coupling. The two-setting
comparison does not assume that a common cancelling curve exists for
both settings.

\subsection{Turning an information bound into a global error bound}
\label{micro:bayes}

A local information bound by itself is not a global risk theorem.
In particular, the reciprocal of information is not automatically an
attainable mean squared error in a phase-estimation problem. Here the
lower bound uses an integrated inequality that remains valid for biased
estimators.

The largest error over a band is at least any weighted average of errors
within it. We can therefore choose convenient weights for a lower-bound
proof, even though the actual experiment makes no probabilistic assumption
about its unknown frequency. Such a weighting is called a prior. We choose
one that goes smoothly to zero at the endpoints so that integration by
parts has no boundary contribution.

Use the frequency-coordinate density
\begin{equation}
 \pi_E(u)=\frac1E\cos^2\!\left(\frac{\pi u}{2E}\right),
 \qquad -E<u<E,
 \qquad
 I(\pi_E)=\int\frac{\pi_E'(u)^2}{\pi_E(u)}\,\dd u
                         =\frac{\pi^2}{E^2}.
 \label{micro:prior}
\end{equation}
This prior is a lower-bound device supported inside the promised band,
not a distribution assumed of the physical frequency. On the nuisance
curve take $\omega=\omega_0+u$. For a fixed protocol with full outcome
likelihood $p(x|u)$, the joint score is
$S=\partial_u\log[p(x|u)\pi_E(u)]$. The density vanishes at both
endpoints, so integration by parts gives
\begin{equation}
 \mathbb E[(\widehat\omega-\omega_0-u)S]=1.
\end{equation}
For a fixed observed outcome, the estimator has no explicit dependence on
the unknown $u$, whereas the true frequency has derivative one. The
derivative of their difference is therefore $-1$, which produces the one
on the right after integration by parts. The conditional likelihood score
has mean zero because differentiating the total outcome probability,
which is one, gives zero. Thus
$\mathbb ES^2=\mathbb E_{\pi_E}J_u+I(\pi_E)$, and another application of
Cauchy--Schwarz yields
\begin{equation}
 \sup_{\omega\in\Omega,\phi}
       \mathbb E(\widehat\omega-\omega)^2
 \ge \frac1{\mathbb E_{\pi_E}J_u+I(\pi_E)}.
 \label{micro:bayesbound}
\end{equation}
No unbiasedness condition appears. This derivation is the scalar
integrated information inequality; the project-specific ingredient is
the curve and the uniform channel bound used within it.

Substitution of Eq.~\eqref{micro:adaptivebound}, together with the
calibration upper bound, proves the explicit fixed-device comparison
\begin{equation}
 \boxed{\quad
 \frac1{nh^2/p_{\min,g}+\pi^2/E^2}
 \le R_g^*(n;\Omega)
 \le\frac1{n\kappa^2}.
 \quad}
 \label{micro:fixedfinite}
\end{equation}
Both denominators have units of time squared, as required. These finite
bounds hold for every integer $n\ge1$.

For a two-setting lower bound, use a different comparison path: keep
$\phi$ fixed. In every selected cell the frequency derivative of the
Hamiltonian remains $Z_S\otimes I_E/2$. Its integrated generator has norm
at most $h/2$. Across $n$ cells, the triangle inequality bounds the
complete purified generator by $nh/2$, including classical selection and
parameter-independent controls. Therefore $J_u\le(nh)^2$. Using the
same prior gives
\begin{equation}
 R_{\{g_a,g_b\}}^*(n;\Omega)
 \ge\frac1{(nh)^2+\pi^2/E^2}.
 \label{micro:menulower}
\end{equation}
The next construction supplies an upper bound $C/n^2$, where $C$ is a
finite constant depending only on the fixed design. An enlarged $C$
handles the finitely many small budgets by using the interval midpoint.
Together these statements prove Theorem~\ref{micro:theorem}.

\subsection{The scope of the adaptive and limiting claims}
\label{micro:limits}

The argument above first covers finite adaptive instruments and their
parameter-independent mixtures. It also extends to the controlled
operational limits included in the model. The global extension rests on
bounded loss, rather than an assumption that Fisher information is
continuous under arbitrary convergence.

The idea can be stated without taking limits of information quantities.
Round each possible estimate to one of finitely many nearby values. The
rounded estimator depends only on which of finitely many outcome bins
occurred. If the probability of each bin converges, its average squared
error converges too. Finally make the rounding finer. Because all estimates
and true frequencies stay in a bounded interval, this last change has a
uniformly small effect on error.

Fix a budget, a comparison curve and its prior. Suppose the full outcome
laws $P_{r,u}$ of finite controllers converge to $P_u$ on one common
measurable output space: for every event and every $u$ on the prior
support, its probability converges. The index $r$ labels an approximating
controller, not an extra resource. Take any measurable estimator $f$ for
the limiting experiment and clip it to $\Omega$. Round its values to a
finite grid in $\Omega$, obtaining $f_\eta$ with
$|f_\eta-f|\le\eta$. The grid inverse images form one fixed finite
partition of the output space. Event-probability convergence gives
pointwise convergence of the risk of $f_\eta$.

All such losses are bounded by $4E^2$. Dominated convergence therefore
passes the same finite-controller Bayes lower bound to $f_\eta$ under
the fixed prior. Finally,
\begin{equation}
 |(f_\eta-\omega)^2-(f-\omega)^2|\le4E\eta
 \qquad (\omega\in\Omega).
\end{equation}
Sending $\eta$ to zero proves the lower bound for $f$ itself. The finite
rounding is classical postprocessing and requires no additional system
contact. This explains exactly which limiting claim is used; agreement
of experiments at a single parameter value would not suffice.

The fixed band remains a substantive physical qualification. A collision
clock can have resonant frequency pairs that give identical complete
boundary-controlled experiments. The theorem promises a suitable
identifiable interval, not uniform consistency at every possible
frequency. Likewise, an instrument placed inside a collision changes the
resource class. Neither an added temporal access resource nor a retuned
compensation is covered by the fixed-device obstruction just proved.


\section{Turning retained records into a global frequency estimate}
\label{phase:section}

The remaining task is constructive.  Two coupling strengths must do more than
produce a large phase derivative at the compensated frequency: they must give
an estimator that works throughout the promised interval, including when the
frequency is initially unknown.  A long interrogation alone cannot do this.
Its phase may pass through many full turns, so a precise phase reading need not
identify a unique frequency.  Moreover, a parity record is random, its
probability depends on frequency, and different records produce different
phase functions.

The construction addresses these difficulties in three steps.  First, an
even-ending record removes the unknown noise direction exactly.  Second, a
gate computed from the record selects a fixed positive fraction of blocks
whose phase functions have a uniformly positive slope.  Third, a sequence of
shorter and longer interrogations identifies the correct phase branch and
refines the frequency estimate.  Every attempted block is charged, whether
or not the gate accepts it.

A \emph{block} is a run of $L$ consecutive cells, followed by a logical
readout and a reset before the next block.  The gate is a calculation on the
recorded outcomes, not an additional physical pulse.  The record determines
a recipe for calculating the phase at any candidate frequency.  Knowing
that recipe is different from knowing its value at the still-unknown true
frequency.

Throughout this section, write $u=\omega-\omega_0$ for residual frequency and
$J=[-E,E]$ for the fixed promised interval.  The variable $v\in J$ will denote
a candidate frequency used in a classical calculation; it is not another
physical setting.  The interval and the hardware are held fixed as the
available number $n$ of collision cells increases.

\subsection{An exact phase attached to an observed path}
\label{phase:identity}

Use the two settings $a$ and $b$ for the couplings $g_a<g_b$.  Start each block
with the even-parity state $|+_L\rangle$, and measure parity at each cell
boundary.  Let $\sigma_i\in\{+1,-1\}$ be the parity after cell $i$, with
$\sigma_0=+1$.  Choose setting $a$ when the preceding parity is even and
setting $b$ when it is odd.  There is no recovery operation after a flip.

The exact single-cell instrument has two normalized branches.  A parity
preserving result applies
\begin{equation}
 D_j(u)=\exp[-i\theta_j(u)Z_S/2]
 \quad\hbox{with probability }q_j(u),
 \label{phase:stay}
\end{equation}
and a parity flip applies $N_\phi$ up to an irrelevant scalar phase, with
probability $p_j(u)=1-q_j(u)$.  These probabilities are independent of the
logical input state and of $\phi$.  The label $j\in\{a,b\}$ is known from
the preceding observed parity.  In particular, the bath need not be
measured: parity distinguishes the two Kraus branches without revealing the
logical amplitudes.

Two elementary identities determine a whole block:
\begin{equation}
 N_\phi^2=I,
 \qquad
 N_\phi D_j(u)=D_j(u)^{-1}N_\phi.
 \label{phase:algebra}
\end{equation}
The second identity follows from $N_\phi Z_SN_\phi=-Z_S$: moving a phase
rotation past a flip reverses its sign.  Move the flip factors together in the
ordered product for an observed record
$\mathcal R=(\sigma_0,\ldots,\sigma_L)$.  If the record ends even, the number
of flips is even, all factors of $N_\phi$ cancel in pairs, and the logical
output is exactly
\begin{align}
 |\psi_{\mathcal R}(u)\rangle
 &=\exp[-i\Phi_{\mathcal R}(u)Z_L/2]|+_L\rangle,\label{phase:state}\\
 \Phi_{\mathcal R}(u)
 &=\sum_{i=1}^{L}\sigma_{i-1}
   \mathbf 1_{\{\sigma_i=\sigma_{i-1}\}}
   \theta_{j(\sigma_{i-1})}(u),
 \qquad j(+1)=a,\quad j(-1)=b.
 \label{phase:recordfunction}
\end{align}
Here $\mathbf 1_A$ equals one when event $A$ occurs and zero otherwise.
$Z_L$ is the logical Pauli operator.  Equation~\eqref{phase:state} is a
statement about each normalized conditional state.  The probability of that
state is still part of the complete experiment.

The unknown direction disappears for every even-ending record and for every
$u$ in the band.  This cancellation is algebraic.  It neither estimates
$\phi$ nor implements an unknown physical operation $N_\phi$ as a control.
The flips in \eqref{phase:algebra} are the recorded noise itself.

For example, consider the five-cell path in Table~\ref{phase:example}.
Its phase is $2\theta_a(u)-\theta_b(u)$ and its probability is
$q_a(u)^2p_a(u)q_b(u)p_b(u)$.  The two flips remove the unknown direction, but
they do not remove the frequency dependence of the record probability.
This small example demonstrates the exact cancellation.  By itself it does
not guarantee that a long record has a useful slope.

\begin{table}[htbp]
\centering
\caption{One even-ending record.  A stay adds a signed signal phase; a flip
changes the sign for subsequent stays.}
\label{phase:example}
\begin{tabular}{ccccc}
\toprule
Cell & Setting & Result & Final parity & Phase contribution\\
\midrule
1 & $a$ & stay & even & $+\theta_a(u)$\\
2 & $a$ & flip & odd & $0$\\
3 & $b$ & stay & odd & $-\theta_b(u)$\\
4 & $b$ & flip & even & $0$\\
5 & $a$ & stay & even & $+\theta_a(u)$\\
\bottomrule
\end{tabular}
\end{table}

\subsection{Why two settings give a positive phase drift}
\label{phase:drift}

At zero residual frequency, define
\begin{equation}
 d_j=\theta_j'(0)=\frac{\tan\chi_j}{g_j},
 \qquad \chi_j=g_jh,
 \qquad
 D_L^0=\sum_{i=1}^{L}\sigma_{i-1}
 \mathbf 1_{\{\sigma_i=\sigma_{i-1}\}}d_{j(\sigma_{i-1})}.
 \label{phase:frozen}
\end{equation}
A prime means differentiation with respect to residual frequency.  Thus each
$d_j$, and $D_L^0$, has units of time, while $\theta_j$ is a dimensionless
angle.  For an even-ending record, $D_L^0=\Phi_{\mathcal R}'(0)$.  The
superscript $0$ means that the known derivatives are evaluated at zero; the
actual record can have been generated at any unknown $u\in J$.
We call each signed contribution a \emph{reward} only as bookkeeping:
positive contributions increase the phase slope, negative ones decrease it.
``Frozen'' means that these coefficients are calculated once at the known
reference point.  It does not mean that the experiment is assumed to be
operating there.

At a fixed actual frequency, the observed parity is a two-state Markov chain
with transition matrix
\begin{equation}
 P(u)=\begin{pmatrix}q_a(u)&p_a(u)\\p_b(u)&q_b(u)\end{pmatrix}.
 \label{phase:matrix}
\end{equation}
Here ``Markov chain'' means that, at this fixed frequency, the probability of
the next parity depends only on the present parity.  The row labels are the
preceding parity, even then odd; the column labels are the subsequent parity
in the same order.  The \emph{stationary fractions} are the long-run
proportions of cells in each parity.  In equilibrium the probability flowing
from even to odd equals the reverse flow: $\pi_+p_a=\pi_-p_b$.
Together with $\pi_++\pi_-=1$, this gives
$\pi_+=p_b/(p_a+p_b)$ and $\pi_-=p_a/(p_a+p_b)$.  For example,
$p_a=0.1$ and $p_b=0.3$ give an even fraction of $3/4$: the system leaves
the odd sector more readily.  This example explains the balance calculation;
the actual phase sensitivity also depends on $d_a,d_b$.
Consequently the stationary
mean frozen reward per cell is
\begin{equation}
 \mu^0(u)=
 \frac{p_b(u)q_a(u)d_a-p_a(u)q_b(u)d_b}{p_a(u)+p_b(u)}.
 \label{phase:driftfunction}
\end{equation}
The positive term counts stays in even parity; the negative term counts
stays in odd parity.  Increasing the coupling while the system is odd
increases its chance of returning to the even sector.  The exact expression
also includes the different conditional phase responses of the two settings.

At $u=0$, let $\mu_0=\mu^0(0)$.  Since
\begin{equation}
 \frac{q_j(0)d_j}{p_j(0)}
 =h\frac{\cot\chi_j}{\chi_j},
 \qquad
 \frac{\mathrm d}{\mathrm d\chi}\frac{\cot\chi}{\chi}
 =-\frac{\chi\csc^2\chi+\cot\chi}{\chi^2}<0,
 \label{phase:positive}
\end{equation}
the strict ordering $0<\chi_a<\chi_b<\pi/4$ implies $\mu_0>0$.
With identical couplings this stationary reward instead vanishes.  This
particular comparison explains the construction's mechanism; the converse
in the preceding section is what excludes all other fixed-coupling
strategies from achieving a better global order.

The fixed band is chosen so that, for every $u\in J$,
\begin{align}
 &0<p_a(u)<p_b(u)<\tfrac12,
 \qquad \mu^0(u)\ge\tfrac34\mu_0,\label{phase:band1}\\
 &\max_{j\in\{a,b\}}\sup_{v\in J}|\theta_j'(v)-d_j|
 \le\tfrac14\mu_0.
 \label{phase:band2}
\end{align}
The functions involved are known and continuous.  Their values at zero give
strict margins for these inequalities, so some interval of positive width
satisfies them.  Once chosen, that interval is fixed.  This argument does
not supply a maximal bandwidth or constants uniform as $g_b-g_a$ tends to
zero.

\subsection{A record gate that certifies an entire phase curve}
\label{phase:gate}

An even endpoint alone is insufficient: signed contributions can largely
cancel on some paths.  Accept a record for phase estimation precisely when
\begin{equation}
 \sigma_L=+1
 \quad\hbox{and}\quad
 D_L^0\ge\frac{\mu_0L}{2}.
 \label{phase:predicate}
\end{equation}
The controller computes this predicate from the observed path and the known
constants $d_a,d_b,\mu_0$.  It does not use the unknown actual frequency or
direction.  A rejected block is retained as a flagged outcome and costs the
same $L$ cells as an accepted block.

Two claims are needed: sufficiently long blocks pass the gate with a
uniformly positive probability, and every accepted record has a useful phase
slope throughout the whole band.  An average slope alone would prove
neither claim.

The next calculation separates the average from its fluctuations.  The mean
score grows proportionally to $L$, whereas its standard deviation grows at
most proportionally to $\sqrt L$.  Thus the relative fluctuation becomes
small in long blocks.  The proof also accounts for the initial even state,
which need not be sampled from the long-run proportions.

\paragraph{Acceptance probability, including the initial transient.}
Set
\begin{equation}
 M=\max(d_a,d_b),\qquad
 \rho=1-\inf_{u\in J}(p_a(u)+p_b(u))<1,
 \qquad
 V=M^2\left(1+\frac{2}{1-\rho}\right).
 \label{phase:varianceconstant}
\end{equation}
For a fixed actual $u$, suppress its argument and write
$\lambda=1-p_a-p_b$.  The band conditions imply $0<\lambda\le\rho<1$.
Starting from even parity, the exact two-state recursion gives
\begin{equation}
 \Pr\{\sigma_i=+1\}=\pi_++\pi_-\lambda^i\ge\tfrac12.
 \label{phase:evenprob}
\end{equation}
Define the conditional mean one-cell rewards
$f_+=q_ad_a$ and $f_-=-q_bd_b$.  Summing their means yields
\begin{equation}
 \mathbb E_u D_L^0
 =L\mu^0(u)+\pi_-(f_+-f_-)
   \sum_{i=0}^{L-1}\lambda^i
 \ge\tfrac34\mu_0 L.
 \label{phase:mean}
\end{equation}
The correction for starting in even parity is nonnegative.  No stationary
initial state has been assumed.

For completeness, the fluctuations also admit a bound uniform over the band.
Write $X_i$ for the $i$th summand of $D_L^0$.  Each $X_i$ lies in $[-M,M]$.
For $i<j$, its conditional prediction at the later cell is
\begin{equation}
 \mathbb E[X_j\mid\sigma_0,\ldots,\sigma_i]
 =(P^{j-i-1}f)(\sigma_i).
 \label{phase:conditionalreward}
\end{equation}
The two possible values on the right differ by at most
$2M\rho^{j-i-1}$, because a two-state chain contracts this difference by
$\lambda$ at each transition.  A random variable in an interval of length
$a$ has variance at most $a^2/4$; applying the covariance inequality to the
two bounded variables therefore gives
\begin{equation}
 |\operatorname{Cov}(X_i,X_j)|\le M^2\rho^{j-i-1},
 \qquad
 \operatorname{Var}(D_L^0)\le VL.
 \label{phase:variance}
\end{equation}
The second inequality follows by summing the geometric covariance bound and
the $L$ individual variances.  It includes the nonstationary start.
Covariance measures whether two rewards fluctuate together.  The geometric
factor says that distant cells remember less about each other's parity, so
their correlations have a finite total contribution per cell.  The bound
on the variance is what turns a positive average into a useful probability
of passing the gate.

An observed reward below $\mu_0L/2$ lies at least $\mu_0L/4$ below its mean.
The variance bound thus implies
\begin{equation}
 \Pr\{D_L^0<\mu_0L/2\}
 \le\frac{16V}{\mu_0^2L}.
 \label{phase:rewardtail}
\end{equation}
Subtract this probability from the lower bound $1/2$ for an even endpoint.
The result is
\begin{equation}
 \Pr\{\mathcal R\hbox{ accepted}\}
 \ge\frac12-\frac{16V}{\mu_0^2L}
 \ge\frac14
 \quad\hbox{for }L\ge L_0,
 \label{phase:acceptance}
\end{equation}
where one sufficient integer is
\begin{equation}
 L_0=\max\left\{1,
 \left\lceil\frac{64V}{\mu_0^2}\right\rceil\right\}.
 \label{phase:minimumlength}
\end{equation}
This subtraction is a union bound; the final parity and the reward need not
be independent.  The value of $L_0$ is a sufficient proof constant and can be
large.  No claim of practical optimality is attached to it.

\paragraph{Sensitivity of each accepted record.}
For any candidate $v\in J$, equation~\eqref{phase:band2} gives the
deterministic bound
\begin{equation}
 |\Phi_{\mathcal R}'(v)-D_L^0|\le\tfrac14\mu_0L.
 \label{phase:derivativecomparison}
\end{equation}
The gate therefore ensures
\begin{equation}
 \alpha L\le\Phi_{\mathcal R}'(v)\le\beta L
 \quad\hbox{for every }v\in J,
 \qquad
 \alpha=\frac{\mu_0}{4},\quad
 \beta=\max_{j\in\{a,b\}}\sup_{v\in J}|\theta_j'(v)|.
 \label{phase:derivativebound}
\end{equation}
Both $\alpha$ and $\beta$ have units of time.  The record certifies a known,
strictly increasing phase function over the whole frequency band, with a
slope proportional to its length.  The probability of receiving such a
certificate is at least $p_0=1/4$, independently of the unknown frequency
and direction.

\subsection{Resolving the phase ambiguity one scale at a time}
\label{phase:multiscale}

The rest of the argument uses only the exact conditional state
\eqref{phase:state}, acceptance probability at least $p_0>0$, and derivative
bounds \eqref{phase:derivativebound}.  It therefore also explains what another
recorded instrument would have to supply for the same estimation method to
apply.  The needed repetitions use fresh bath cells and a full reset to
$|+_L\rangle$ between blocks.

Before specifying the proof constants, the procedure is:
\begin{enumerate}
\item Use short, paid trials to find a coarse frequency interval.
\item Double the block length and predict each accepted record's phase at
the midpoint of that interval.
\item Measure relative to that prediction, compare the outcomes with the
recorded response functions, and obtain a refined estimate.
\item Intersect a narrower window around the estimate with the old interval,
then repeat at the next length.
\end{enumerate}
Long blocks give sharper readings, while the preceding shorter blocks tell
us which phase turn we are on.  Intersection keeps all later estimates
inside the earlier interval, even if a later refinement fails.  The constants
below are one sufficient way to make both the failure probabilities and the
total cost precise.

The paid one-cell calibration experiment supplies an estimator
$\widehat u_N\in J$ with
\begin{equation}
 \sup_{u\in J,\phi}
 \Pr\{|\widehat u_N-u|\ge\delta\}
 \le2\exp[-N\kappa_0^2\delta^2/2],
 \qquad \delta>0.
 \label{phase:coarsebound}
\end{equation}
Here $\kappa_0=\kappa>0$ is the uniform derivative lower bound for the one-cell
calibration response, with units of time.  A fixed initial localization
accuracy will suffice, but its failure probability must become smaller as
the final experiment becomes longer.

Choose a final scale $K\ge1$ and define
\begin{equation}
 s=\frac{\pi}{8\beta},\qquad
 L_k=2^kL_0,\qquad
 \gamma=\min\left\{\frac{p_0^2}{2},
                   \frac{p_0\alpha^2s^2}{8}\right\}>0.
 \label{phase:scheduleconstants}
\end{equation}
$L_k$ counts cells, $s$ has units of frequency, and $\gamma$ is
dimensionless.  The symbol $\pi$ in this equation is the usual circle
constant.  All logarithms below are natural logarithms.

First use
\begin{equation}
 N_0(K)=\left\lceil
 \frac{2L_0^2}{\kappa_0^2s^2}
       (K\log 8+\log 2)\right\rceil
 \label{phase:calibrationcost}
\end{equation}
one-cell calibration trials and form
\begin{equation}
 C_0=J\cap[\widehat u_0-s/L_0,\widehat u_0+s/L_0].
 \label{phase:initialinterval}
\end{equation}
The estimate $\widehat u_0$ is clipped to $J$, so $C_0$ is always nonempty.
Its diameter is at most $2s/L_0$, and
$\Pr\{u\notin C_0\}\le8^{-K}$.  This initial step costs only a number of
cells proportional to $K$, despite the increasingly small failure
probability, because its target width stays fixed.

At scale $k=1,\ldots,K$, let $c$ be the midpoint of the current interval
$C_{k-1}$.  Hold that center and the block policy fixed throughout this
stage, and attempt
\begin{equation}
 r_k=\left\lceil
 \frac{(K-k)\log8+\log12}{\gamma}
 \right\rceil
 \label{phase:repetitioncounts}
\end{equation}
independent blocks of length $L_k$.  Earlier stages may determine $c$.
Earlier terminal measurement outcomes within this same stage do not change
the later block policies.  This restriction is part of the explicit
attaining strategy and will make its conditional independence transparent.

For each accepted record $\mathcal R_i$, measure the logical equatorial
observable
\begin{equation}
 O_i=\cos(\Phi_{\mathcal R_i}(c))Y_L
       -\sin(\Phi_{\mathcal R_i}(c))X_L.
 \label{phase:observable}
\end{equation}
Here $X_L,Y_L$ are the logical Pauli operators.  The outcome
$Y_i\in\{-1,+1\}$ is a classical number; the subscript distinguishes it from
the operator $Y_L$.  The conditional mean is exactly
\begin{equation}
 \mathbb E_u[Y_i\mid\mathcal R_i]
 =\sin(\Phi_{\mathcal R_i}(u)-\Phi_{\mathcal R_i}(c)).
 \label{phase:meanreadout}
\end{equation}
This centers the readout at the phase predicted by the current interval
center.  It changes a terminal measurement basis only.  The compensation,
cell duration and interaction Hamiltonian remain fixed.

Let $G$ be the number of accepted blocks in the stage.  If
$G<p_0r_k/2$, take $\widetilde u_k=c$.  Otherwise, compute from the actual
accepted records the known function
\begin{equation}
 H_k(v)=\frac1G\sum_{i:\,\mathcal R_i\text{ accepted}}
 \sin(\Phi_{\mathcal R_i}(v)-\Phi_{\mathcal R_i}(c)),
 \qquad v\in C_{k-1}.
 \label{phase:responsecurve}
\end{equation}
Average the corresponding measured $Y_i$, clip that average to the image
interval $H_k(C_{k-1})$, and invert $H_k$ to obtain $\widetilde u_k$.  The
clipping prevents a fluctuation from requesting an inverse outside the
current candidate interval.  Finally set
\begin{equation}
 C_k=C_{k-1}\cap
 [\widetilde u_k-s/L_k,\widetilde u_k+s/L_k].
 \label{phase:nesting}
\end{equation}
Every interval remains nonempty and nested, even when a statistical failure
has occurred: $\widetilde u_k$ belongs to the preceding interval and to the
new window.  If an interval has already collapsed to a point, that point is
the estimate and no inversion is needed.

Why is the inverse unique?  The preceding interval has diameter at most
$2s/L_{k-1}$, so every $v\in C_{k-1}$ obeys $|v-c|\le s/L_{k-1}$.
For every accepted record,
\begin{equation}
 |\Phi_{\mathcal R_i}(v)-\Phi_{\mathcal R_i}(c)|
 \le\beta L_k\frac{s}{L_{k-1}}=\frac\pi4.
 \label{phase:noalias}
\end{equation}
All candidate phases therefore lie on a single increasing branch of the
sine function.  Combining its derivative with
\eqref{phase:derivativebound} gives
\begin{equation}
 H_k'(v)\ge\frac{\alpha L_k}{\sqrt2}.
 \label{phase:inversebound}
\end{equation}
The full phase $\Phi_{\mathcal R_i}(v)$ may have wound around the circle many
times.  What is small is its difference across the candidate interval, made
small by the earlier paid localization.  This is how the construction
resolves phase ambiguity instead of assuming it away.

\subsection{Conditioning on the record, without freezing its law}
\label{phase:conditioning}

For example, if two records are accepted, their prediction is the average
of their two separate sine responses,
\begin{equation}
 H_k(v)=\tfrac12\bigl[
 \sin(\Phi_1(v)-\Phi_1(c))+
 \sin(\Phi_2(v)-\Phi_2(c))\bigr].
\end{equation}
The record fixes which functions appear here.  After fixing it, the remaining
randomness is in the final $+1/-1$ measurements.  A concentration bound says
that their measured average is unlikely to differ greatly from this predicted
average.  Because the bound holds for every accepted collection of functions,
it can then be averaged over their actual frequency-dependent occurrence
probabilities.  No average phase is substituted inside the sine.
The two-record example illustrates the prediction rule, not a sufficient
number of repetitions for a precise estimate.

The law of $\mathcal R_i$ depends on $u$ through $p_a(u)$ and $p_b(u)$.
This dependence does not invalidate \eqref{phase:responsecurve}.  For the
fixed unknown value $u$, first condition on all earlier stages.  The fresh
blocks have independent acceptance indicators and each succeeds with
probability at least $p_0$.  Bounded-variable concentration gives
\begin{equation}
 \Pr\{G<p_0r_k/2\}\le\exp[-p_0^2r_k/2].
 \label{phase:counttail}
\end{equation}
Next condition on all the records of this stage.  Because the policy was
held fixed and the blocks used fresh bath cells, the terminal outcomes of
the accepted blocks are conditionally independent.  They need not have the
same conditional mean; their average conditional mean is exactly $H_k(u)$.

When $u\in C_{k-1}$, clipping and
\eqref{phase:inversebound} imply that an estimation error exceeding
$s/L_k$ requires the sample mean to differ from $H_k(u)$ by more than
$\alpha s/\sqrt2$.  Independent variables in $[-1,1]$ obey the tail bound
$2\exp(-Gt^2/2)$ for a deviation $t$ of their average from its mean.
Consequently, on $G\ge p_0r_k/2$,
\begin{equation}
 \Pr\{|\widetilde u_k-u|>s/L_k\mid\mathcal R_1,\ldots,\mathcal R_{r_k}\}
 \le2\exp[-G\alpha^2s^2/4]
 \le2\exp[-p_0r_k\alpha^2s^2/8].
 \label{phase:readouttail}
\end{equation}
Average this inequality over the actual record law and add the low-count
event.  Uniformly over earlier observations for which $u\in C_{k-1}$,
\begin{equation}
 \Pr\{u\notin C_k\mid\hbox{earlier observations}\}
 \le3e^{-\gamma r_k}
 \le\frac14\,8^{-(K-k)}.
 \label{phase:stagefailure}
\end{equation}
No derivative of the record distribution was taken, and no record
probability was replaced by its compensated value.  In particular,
$H_k$ is the average of the phase responses for the records actually seen,
not the expectation of a phase under a guessed parameter-independent law.
Classical information in record probabilities could also be used, but is
not needed for this attainable bound.

\subsection{Worst-case error with a hard cell budget}
\label{phase:riskcost}

Return the midpoint $\widehat u$ of $C_K$.  If the coarse localization and
every refinement succeed, then
$|\widehat u-u|\le s/L_K$.  If the first failed refinement is stage $k$,
the true frequency and every later estimate still belong to $C_{k-1}$.
Thus their final separation is at most its diameter,
$2s/L_{k-1}=4s/L_k$.  If the coarse step fails, both the true value and the
final estimate still lie in $J$, giving the bound $2E$.

These three alternatives give an unconditional, worst-case mean-square
error bound:
\begin{align}
 \sup_{u\in J,\phi}\mathbb E[(\widehat u-u)^2]
 &\le\frac{s^2}{L_K^2}+4E^2\,8^{-K}
 +\sum_{k=1}^{K}\frac{16s^2}{L_k^2}\,
                   \frac14\,8^{-(K-k)}\notag\\
 &\le\frac{9s^2+4E^2L_0^2}{L_K^2}.
 \label{phase:globalrisk}
\end{align}
The last step is a geometric sum: after multiplication by $L_K^2$, the
refinement terms are bounded by
$4s^2\sum_{j=0}^{K-1}2^{-j}\le8s^2$.  The coarse term uses
$4^K8^{-K}\le1$.  This proof includes failed branches instead of reporting
only the precision conditional on success.

The total number of cells is the deterministic number
\begin{equation}
 B_K=N_0(K)+\sum_{k=1}^{K}r_kL_k.
 \label{phase:budget}
\end{equation}
The repetition count is largest at the short scales, where a failure could
cause the largest final error.  Long, costly blocks require fewer
repetitions because a late failure can move the estimate only within an
already narrow interval.  In particular,
\begin{equation}
 \sum_{k=1}^{K}L_k\le2L_K,
 \qquad
 \sum_{k=1}^{K}(K-k)L_k\le2L_K.
 \label{phase:costsums}
\end{equation}
Since $N_0(K)=O(K+1)$ and $K+1\le2^K$ for $K\ge1$, there is a fixed
finite constant $C_*$, depending only on the hardware and chosen band, with
\begin{equation}
 L_K\le B_K\le C_*L_K.
 \label{phase:linearcost}
\end{equation}
There is no extra logarithmic factor in the total cost.  The decreasing
repetition schedule is essential to this conclusion; requiring the same
small failure probability at every scale would waste resources.

For a prescribed hard budget $n$, select the largest feasible $K$.
Then $n<B_{K+1}\le2C_*L_K$, so \eqref{phase:globalrisk} becomes
\begin{equation}
 \sup_{\omega\in[\omega_0-E,\omega_0+E],\phi}
 \mathbb E[(\widehat\omega-\omega)^2]=O(n^{-2}),
 \qquad \widehat\omega=\omega_0+\widehat u.
 \label{phase:attainment}
\end{equation}
The finitely many budgets too small for the first scale use the interval
midpoint; increasing the fixed upper-bound constant covers them.  The
matching inverse-square lower bound in the preceding section therefore
establishes the two-coupling minimax order.  This construction does not
determine the optimal global risk coefficient.

\paragraph{Contacts and duration.}
Each block has its prescribed length even when its record fails the gate.
At a cell boundary, parity extraction, a completed block's logical readout,
the next preparation and a coupling choice can be consecutive parts of one
allowed instrument.  Scheduling all used cells consecutively within the
$n$-cell allocation therefore uses at most $n-1$ nonterminal boundary
contacts, with the final readout at $nh$ free under the stated convention.
If a complete estimate is produced earlier, that readout is a counted
contact and the classical result is retained through the remaining charged
cells.  The physical duration remains $nh$.

Calibration, failed gates, resets, repetitions and premature terminal
readouts have all appeared in the same accounting.  At no point does a
classical frequency estimate alter the physical compensation.  The
quadratic precision comes from exact noise-direction cancellation and
usable recorded phase functions, combined with paid ambiguity resolution.
An unrecorded mixture, a phase known only through its derivative at zero,
or an unknown timing perturbation would require a new argument.


\section{The continuous-time counterpart}
\label{markov:continuous}

The finite-collision result has a related continuous-time statement. Its
importance is conceptual: the obstruction caused by an unknown error direction
does not depend on a weak-collision approximation. The continuous model also
admits an especially transparent way to convert noise into an exact, useful
interrogation. It is a separate model, with ideal continuous syndrome monitoring.
We do not obtain its theorems by silently replacing the finite collision clock
by arbitrarily fast access inside a collision.

\subsection{What changes, and what remains the same}

The sensor is a qubit, the auxiliary memory is noiseless, and the target
frequency is $\omega$. A noise rate $\gamma_j$ has a direct operational
meaning: during a short time $\mathrm dt$, a flip occurs with probability
$\gamma_j\mathrm dt$ up to terms smaller than $\mathrm dt$. Between flips,
the signal rotates the state. At a fixed setting, the probability of a flip
in the next short interval does not depend on how long the system has already
waited. This is the memoryless, or Markovian, assumption in this model.
Averaging the flip and no-flip alternatives gives the equation below.
In a selected setting $j$, the reduced sensor state obeys
\begin{equation}
 \frac{\mathrm d\rho}{\mathrm dt}
 =-i\left[\frac{s_j\omega}{2}Z,\rho\right]
   +\gamma_j(N_\phi\rho N_\phi-\rho),
 \qquad N_\phi=\cos\phi\,X+\sin\phi\,Y.
 \label{markov:master}
\end{equation}
Here $\rho$ is a density matrix, $[A,B]=AB-BA$, the known real number $s_j$
is a dimensionless signal multiplier, and the known rate $\gamma_j>0$ has units
of inverse time. The angle $\phi$ is unknown, continuous, fixed throughout the
experiment, and shared by every setting. Its physical period in this
dissipator is $\pi$, because changing the sign of $N_\phi$ leaves
\eqref{markov:master} unchanged. In the main comparison $s_j=1$: classical rate selection allows
$\gamma(t)\in[a,b]$ with fixed $0<a<b$, while preserving the signal
Hamiltonian. The fixed-device benchmark uses one specified rate $\gamma>0$.
The endpoints $\{a,b\}$ suffice for the attaining protocols below; the
converse covers every rate in the interval. The more general known
coefficients $s_j$ are retained for the finite-setting extension.

The resource $T$ is a hard physical exposure time for one sensor. Preparation,
parameter-independent sensor--memory controls, monitoring, recovery pulses,
classical feedback and protected storage are ideal. These idealizations are
common to both sides of the comparison. They do not provide direct bath access.
Rate selection means a classical choice of the scalar multiplying the same
unknown noise operator. It does not mean a coherent superposition of baths or
coherent access to $N_\phi$ as an extra Hamiltonian.

Use the two-qubit code $|0_L\rangle=|00\rangle$ and
$|1_L\rangle=|11\rangle$, with the sensor first. The signal preserves parity;
an error $N_\phi$ changes it. Ideal parity monitoring therefore identifies
errors without measuring the encoded amplitudes. The jump operator is
$L_j=\sqrt{\gamma_j}N_\phi$, so $L_j^\dagger L_j=\gamma_j I$.
Consequently the jump rate is independent of the logical state, $\omega$ and
$\phi$. For a constant setting the observed jump count is Poisson with mean
$\gamma_j t$: explicitly,
$\Pr\{N=k\}=e^{-\gamma_jt}(\gamma_jt)^k/k!$ for
$k=0,1,\ldots$. Here $k!$ is the product of the positive integers up to $k$,
with $0!=1$. Both the mean and variance of this count equal $\gamma_jt$.
Thus the typical count is proportional to elapsed time, with fluctuations
of order its square root. Between jumps, normalized evolution is just the signal rotation.
These facts make the complete measurement record tractable; none permits
discarding the time spent on an unsuccessful run.

\subsection{Why a large frequency entry in an information matrix is insufficient}

For an outcome distribution $p(x\mid\omega,\phi)$, its Fisher matrix has entries
$J_{uv}=\mathbb E[(\partial_u\log p)(\partial_v\log p)]$,
where $u,v\in\{\omega,\phi\}$. A nuisance direction can imitate a target
change. The information that remains after allowing that imitation is
\begin{equation}
 \mathcal I_\omega(J)
 =\inf_{c\in\mathbb R}(1,c)J(1,c)^{\mathsf T}
 =J_{\omega\omega}-\frac{J_{\omega\phi}^2}{J_{\phi\phi}}
 \quad\hbox{when }J_{\phi\phi}>0.
 \label{markov:efficient}
\end{equation}
The infimum definition also handles a zero or singular nuisance block. The
superscript $\mathsf T$ denotes transpose. Frequency information has units of
time squared. Its reciprocal bounds the variance of a regular locally unbiased
frequency estimator, with zero first-order response to $\phi$. This is a local
statement about a specified parameter point. It is not yet a global risk bound
for an initially unknown frequency.

An elementary exact example shows what is at stake. Apply the fixed sensor
recovery $X$ after every detected error. Since $XN_\phi=e^{i\phi Z}$, an
interrogation of duration $t$ and observed count $N$ has phase
\begin{equation}
 \Theta=\omega t-2\phi N,\qquad N\sim\mathrm{Poisson}(\gamma t).
 \label{markov:recoveredphase}
\end{equation}
For a typical long run, $N$ is near $\gamma t$, so the phase's leading
time-proportional part is $(\omega-2\gamma\phi)t$. Count fluctuations supply
the remaining distinction. A change in frequency can therefore be largely
imitated by a change in the unknown direction. Two rates produce different
combinations of the same two unknowns; this is the statistical reason that
rate diversity can help.
Prepare the logical $X$ eigenstate and measure at a uniformly random,
recorded equatorial angle $\vartheta$. The two probabilities are
$[1\pm\cos(\Theta-\vartheta)]/2$. This measurement has unit information about
$\Theta$ except at isolated angles, which have probability zero under the
recorded continuous randomization. The count likelihood itself contains no
information about $\omega$ or $\phi$.

More generally, a recorded phase $\omega D-2\phi M$, with a
parameter-independent distribution of the real variables $(D,M)$, gives
\begin{equation}
 J=\begin{pmatrix}
 \mathbb E D^2&-2\mathbb E(DM)\\
 -2\mathbb E(DM)&4\mathbb E M^2
 \end{pmatrix}.
 \label{markov:phasefim}
\end{equation}
For \eqref{markov:recoveredphase}, $D=t$,
$\mathbb E N=\gamma t$ and $\mathbb E N^2=(\gamma t)^2+\gamma t$.
Substitution in \eqref{markov:efficient} gives
\begin{equation}
 \mathcal I_\omega(t)=\frac{t^2}{1+\gamma t}.
 \label{markov:fixedattain}
\end{equation}
The entry $J_{\omega\omega}=t^2$ is large, but its leading part is
indistinguishable from an angle change. Only linear growth remains at long
times. All counts and readout labels are retained in this calculation.

\subsection{A converse for every adaptive control strategy}

The limitation is not specific to that recovery. Restrict the two unknown
parameters to a mathematical comparison path
$\omega(u)=\omega_c+u$, $\phi(u)=\phi_c+cu$. The real parameter $u$
has frequency units, and $c$ has time units. The controller is held fixed along
this path. Outcome probabilities may change; the control rule does not acquire
knowledge of the unknown $u$.

We use a general information-growth inequality for a Markovian channel with
arbitrary parameter-independent controls. If $F_u(t)$ is the output quantum
Fisher information along the path, including the classical record, then
\begin{align}
 \frac{\mathrm dF_u}{\mathrm dt}
 &\le 4\bigl(\|A\|+\|B\|\sqrt{F_u}\bigr),\label{markov:growth}\\
 A&=(\partial_uL)^\dagger\partial_uL,\qquad
 B=\partial_uH-\frac{i}{2}
 \bigl((\partial_uL)^\dagger L-L^\dagger\partial_uL\bigr).
 \nonumber
\end{align}
The operator norm $\|A\|$ is the largest factor by which $A$ can stretch
a vector. Quantum Fisher
information bounds the classical Fisher information of every final
measurement. The inequality is an information-theoretic ingredient; the
following choice of nuisance path is its application to this physical model.
The two terms have different roles. The $A$ term allows information to
increase at most at a bounded rate, giving a linear ceiling if $B=0$.
The term containing $B$ permits
coherent accumulation: because its size increases with $\sqrt{F_u}$, it can
support quadratic growth. We look for a nuisance path that cancels this
second contribution. One rate admits such a cancellation; two different
rates generally do not admit the same cancelling path.
For adaptive setting selection, include the controller and recorded branches
as orthogonal blocks and take the largest norms over the permitted settings.
These blocks simply keep mutually distinguishable classical records in
separate labels of the mathematical description.
Measurements, resets and feedback remain parameter-independent channels. This
is why the bound includes learning the angle during the experiment.

The Pauli identities
$(\partial_\phi N_\phi)N_\phi=-iZ$ and
$N_\phi(\partial_\phi N_\phi)=iZ$ give, in setting $j$,
\begin{equation}
 A_j=c^2\gamma_j I,\qquad
 B_j=\left(\frac{s_j}{2}-c\gamma_j\right)Z.
 \label{markov:ab}
\end{equation}
For one unit-strength setting, $c=1/(2\gamma)$ cancels $B$ exactly.
Starting from a parameter-independent preparation, $F_u(0)=0$, so
$F_u(T)\le T/\gamma$. The path $c=0$ separately gives $F_u(T)\le T^2$.
Together with \eqref{markov:fixedattain}, this proves
\begin{equation}
 \mathcal I^*_{\omega,\mathrm{fixed}}(T)\le\min\{T^2,T/\gamma\},
 \qquad
 \lim_{T\to\infty}\frac{\mathcal I^*_{\omega,\mathrm{fixed}}(T)}{T}
 =\frac1\gamma.
 \label{markov:fixedlocal}
\end{equation}
The star denotes optimization over all permitted controls and measurements.
The mathematical path does not assert that noise physically encodes frequency;
it selects difficult pairs inside the allowed two-parameter family.

For rates selected from $[a,b]$, define the dimensionless contrast
\begin{equation}
 r=\frac{b-a}{b+a}\in(0,1).
 \label{markov:r}
\end{equation}
Choosing $c=1/(a+b)$ leaves $\|B_\gamma\|\le r/2$ for every $\gamma\in[a,b]$ and
$\|A_\gamma\|\le b/(a+b)^2$. Thus
\begin{equation}
 \frac{\mathrm d F_u}{\mathrm dt}\le C+2r\sqrt{F_u},\qquad
 C=\frac{4b}{(a+b)^2}.
 \label{markov:ode}
\end{equation}
The derivative $\mathrm dF_u/\mathrm dt$ is with respect to physical time. Integrating the scalar comparison
equation gives $\sqrt{F_u(t)}\le rt+O(t_c\log(1+t/t_c))$ at large $t$,
where $t_c=1/(a+b)$ and the rates are fixed; hence the leading information
coefficient is at most $r^2$.
The comparison integral is
$t=y/r-[C/(2r^2)]\log(1+2ry/C)$, where $y=\sqrt{F_u}$ and the logarithm
is natural. This also explains why the term linear in time cannot change the
leading quadratic coefficient.

\subsection{Rate diversity can remove the leading statistical ambiguity}

Two simple constructions attain that coefficient without adapting the bath
rate to the syndrome. They still use syndrome-triggered recovery pulses.

First, use rate $a$ for time $t_a$, rate $b$ for time $t_b$, with
$t_a+t_b=T$. Insert and later undo a logical $X$ pulse so that the two block
phases subtract. With independent Poisson counts $N_a,N_b$,
\begin{equation}
 \Theta=\omega D-2\phi M,\quad
 D=t_a-t_b,\quad M=N_a-N_b,\quad
 \mu=at_a-bt_b,\quad V=at_a+bt_b.
\end{equation}
Since $\mathbb E M=\mu$ and $\mathbb E M^2=V+\mu^2$,
\begin{equation}
 \mathcal I_{\rm echo}=\frac{D^2V}{V+\mu^2}.
 \label{markov:echo}
\end{equation}
The choice $t_a=bT/(a+b)$, $t_b=aT/(a+b)$ makes $\mu=0$ and $D=rT$.
It achieves $\mathcal I_{\rm echo}=r^2T^2$ exactly at every $T$.
The nuisance remains in individual output states, but its score is statistically
orthogonal to the frequency score after the complete records are pooled.

This design has a specific calibration sensitivity. Put $t_a=fT$ and define
$q=af-b(1-f)$, $v=af+b(1-f)$. Then
\begin{equation}
 \mathcal I_{\rm echo}
 =\frac{(2f-1)^2T^2}{1+(q^2/v)T}.
 \label{markov:imbalance}
\end{equation}
A fixed nonzero imbalance $q$ eventually restores linear information growth.
The dimensionless quantity $q^2T/v$ must stay bounded to retain a nonzero
quadratic coefficient. This is a restriction of this particular balanced echo,
not a universal accuracy requirement for rate-diverse sensing.

Second, choose rate $a$ with probability $p$ and $b$ otherwise at the start of
an interrogation. Retain the choice and use it throughout duration $T$.
Define $\bar\gamma=pa+(1-p)b$, $g_2=pa^2+(1-p)b^2$ and
$\sigma_\gamma^2=p(1-p)(b-a)^2$. Pooling the full Fisher matrices before
eliminating their \emph{shared} nuisance gives
\begin{equation}
 \mathcal I_p(T)
 =T^2\frac{\sigma_\gamma^2+\bar\gamma/T}{g_2+\bar\gamma/T},
 \qquad
 \lim_{T\to\infty}\frac{\mathcal I_p(T)}{T^2}
 =\frac{p(1-p)(b-a)^2}{pa^2+(1-p)b^2}.
 \label{markov:mixture}
\end{equation}
The maximizing probability is $p=b/(a+b)$ and the maximum is $r^2$.
These endpoint-only experiments attain the interval converse, so
\begin{equation}
 \lim_{T\to\infty}\frac{\mathcal I^*_{\omega,[a,b]}(T)}{T^2}=r^2.
\end{equation}
A fixed error in $p$ changes this coefficient continuously; it does not incur
the time-amplified penalty in \eqref{markov:imbalance}. A fair choice $p=1/2$
needs no numerical rate calibration in its design and retains the fraction
$(a+b)^2/[2(a^2+b^2)]>1/2$ of the optimal local coefficient. For example,
$b=3a$ gives $r^2=1/4$; the fair design retains four fifths of that coefficient.

Unknown numerical rates can be included as additional nuisance parameters in
these likelihoods. With fixed durations and selection probabilities, the phase
likelihood conditional on counts depends on $(\omega,\phi)$, while the rate
score comes from the count likelihood. Their cross-information is zero because
the conditional phase score has mean zero. The displayed efficient information
therefore survives elimination of the rate parameters. This fact is not free
calibration: inference uses the recorded counts, and a precisely balanced echo
still needs its design ratio. Nor does one labelled interrogation establish
global estimation at an unknown frequency. Repeating it $N_{\rm rep}$ times costs
$N_{\rm rep}T$.

There is a structural extension. For finitely many selectable pairs
$(s_j,\gamma_j)$, the optimal leading local coefficient is $v_*^2$, where
\begin{equation}
 v_* =\max_{i,j}
 \frac{|s_i\gamma_j-s_j\gamma_i|}{\gamma_i+\gamma_j}.
 \label{markov:general}
\end{equation}
Indeed the converse minimizes $\max_j|s_j-2c\gamma_j|$. Requiring this
maximum to be at most $v$ says that the intervals
$[(s_j-v)/(2\gamma_j),(s_j+v)/(2\gamma_j)]$ have a common intersection.
Their endpoint inequalities give precisely \eqref{markov:general}. A recorded
random choice of a maximizing pair, with probabilities proportional to the
other member's rate, attains it by \eqref{markov:phasefim}. If all ratios
$s_j/\gamma_j$ coincide, one nuisance path cancels every $B_j$, and the
stronger bound $\mathcal I^*_\omega(T)\le T\max_j(s_j^2/\gamma_j)$ holds.
Speeding up signal and noise in the same proportion therefore does not remove
the obstruction. If each setting has an independently unknown direction, one
can instead choose a separate cancelling nuisance slope for each setting; the
same linear bound returns. A stable shared direction is essential.

\section{Exact stopped sensing and a globally unknown frequency}
\label{markov:global}

Local information identifies a promising resource, but it does not resolve
the multiple fringes of a long phase measurement. A different construction
supplies a direction-independent channel on every completed run and makes
the global precision theorem possible.

\subsection{Use errors in pairs, and stop a signed clock}

Again monitor parity, but apply no recovery after an error. Let $N(t)$ be the
detected jump count and define the signed interrogation time
\begin{equation}
 D(t)=\int_0^t(-1)^{N(u)}\,\mathrm du.
 \label{markov:signed}
\end{equation}
It advances at unit speed in even parity and retreats at unit speed in odd
parity. For example, a trajectory spending 0.9 time units even, 0.25 odd,
and another 0.85 even has physical duration 2 but signed duration
$0.9-0.25+0.85=1.5$. The negative interval subtracts signal phase rather
than physical cost. Since $N_\phi Z=-ZN_\phi$, moving errors through the intervening
signal rotations gives the exact trajectory unitary
\begin{equation}
 U(t)=N_\phi^{N(t)}\exp[-i\omega D(t)Z/2].
 \label{markov:trajectory}
\end{equation}
Use rate $a$ when parity is even and $b$ when it is odd. The positive state
lasts longer on average, so the signed clock has positive long-time drift $r$.
For a prescribed $d>0$, stop at its first arrival at $d$:
\begin{equation}
 \tau_d=\inf\{t\ge0:D(t)=d\}.
\end{equation}
An upward arrival is in even parity, so the number of errors is even and
$N_\phi^{N(\tau_d)}=I$. The completed logical operation is exactly
\begin{equation}
 \rho_L\longmapsto e^{-i\omega dZ_L/2}\rho_L e^{i\omega dZ_L/2}.
 \label{markov:stoppedchannel}
\end{equation}
It is correct for every logical input, frequency and direction. The stopping
rule requires neither $\phi$ nor the numerical rates. Rates determine its
cost, and rate knowledge enters the performance guarantees below. The
stopping-time record is independent of $\omega$ and $\phi$.

\subsection{The cost of waiting is part of the result}

The time needed to reach the prescribed signed duration is random. Its mean
alone cannot enforce a deadline: rare runs might last much longer. We therefore
describe its distribution through weighted averages of the form
$\mathbb E e^{-q\tau_d}$. For $q>0$, this gives less weight to longer runs.
As $q$ varies, these averages form the \emph{Laplace transform} of the waiting
time. Its first derivative at zero, with a minus sign, gives the mean; the
second derivative of its logarithm gives the variance. Positive exponential
weights $e^{\theta\tau_d}$ instead emphasize unusually long runs and will
control deadline failures.

Write $\delta=(b-a)/2$. For a Laplace variable $q\ge0$, with units of
inverse time, the exact first-passage transform is
\begin{equation}
 \mathbb E e^{-q\tau_d}=e^{-d k(q)},\qquad
 k(q)=\sqrt{q^2+(a+b)q+\delta^2}-\delta.
 \label{markov:laplace}
\end{equation}
To see where it comes from, let $f_\pm(x)$ be the transform of the remaining
time to $d$ when the signed clock is at $x<d$ in even or odd parity.
Conditioning on the next infinitesimal step gives
\begin{align}
 \partial_x f_+ +a(f_- -f_+)&=qf_+,\nonumber\\
 -\partial_x f_- +b(f_+-f_-)&=qf_-.
 \label{markov:backward}
\end{align}
The form $f_\pm(x)=v_\pm e^{-k(d-x)}$ gives
$k^2+(b-a)k-q(q+a+b)=0$. The outgoing boundary is $v_+=1$;
the incoming value $f_-(d^-)$ is not assigned that same boundary condition.
The root bounded towards $x=-\infty$ gives \eqref{markov:laplace} for the
initial state $x=0$, even parity.

At $q=0$ the transform equals one, establishing almost-sure completion,
meaning completion with probability one.
Differentiating its logarithm at zero gives
\begin{equation}
 \mathbb E\tau_d=\frac d r,\qquad
 \operatorname{Var}(\tau_d)=\frac{8ab\,d}{(b-a)^3}.
 \label{markov:moments}
\end{equation}
The transform also has a finite positive exponential-moment domain:
\begin{equation}
 \mathbb E e^{\theta\tau_d}=e^{-d k(-\theta)},\qquad
 0<\theta<\theta_c,\qquad
 \theta_c=\frac{(\sqrt b-\sqrt a)^2}{2}.
 \label{markov:mgf}
\end{equation}
Finiteness can be established before continuation of the transform. For a
real tilt $\eta$ with inverse-time units, exponential moments of $D(t)$ are
governed by the two-state matrix
\begin{equation}
 Q_\eta=\begin{pmatrix}\eta-a&a\\b&-\eta-b\end{pmatrix},\qquad
 \lambda(\eta)=-\frac{a+b}{2}+\sqrt{(\eta+\delta)^2+ab}.
\end{equation}
The quantity $\lambda(\eta)$ is its largest eigenvalue. At $\eta=-\delta<0$ it
equals $-\theta_c$. For fixed rates, the corresponding
matrix exponential bounds $\mathbb E e^{-\delta D(t)}$ by a constant times
$e^{-\theta_ct}$. Since $\tau_d>t$ implies $D(t)<d$, the exponential
Markov inequality bounds its probability by a constant times
$e^{\delta d-\theta_ct}$. Thus positive exponential moments of
$\tau_d$ are finite throughout the domain in \eqref{markov:mgf}. Continuing
\eqref{markov:laplace} within that domain yields the stated formula.

The strict inequality $a<b$ matters. If $a=b=\gamma$, the transform becomes
$\exp[-d\sqrt{q(q+2\gamma)}]$. Completion remains almost sure, but its mean
is infinite: increasingly rare but increasingly long waits make the average
diverge. An eventual noiseless output alone therefore supplies no
finite-cost precision advantage.

Under a deadline $T$, retain both success and failure flags. Transfer completed
logical outputs to protected memory; on failure assign a parameter-independent
output. The experiment has local information
$\Pr(\tau_d\le T)d^2$. For instance, choose $t_c=1/(a+b)$,
$\Delta_T=t_c(T/t_c)^{2/3}$ and $d=r(T-\Delta_T)$ at large $T$. The mean duration falls short of
$T$ by the stated margin, while its variance is $O(T)$. The failure probability
is $O(T^{-1/3})$ by the variance bound, and $\Pr(\tau_d\le T)d^2/T^2\to r^2$.
Thus stopped sensing also attains the optimal local coefficient with every
failed run charged. This margin is an explicit construction, not a claim about
the best finite-time stopping design.

\subsection{The same global risk criterion on both sides}

Fix a closed frequency interval $\Omega$ of width $W>0$, independent of $T$.
For the permitted control menu $\mathcal A$, define
\begin{equation}
 R_{\mathcal A}^*(T;\Omega)
 =\inf_{\mathcal P,\widehat\omega}
 \sup_{\omega\in\Omega,\,\phi}
 \mathbb E_{\omega,\phi}^{\mathcal P}
       (\widehat\omega-\omega)^2.
 \label{markov:risk}
\end{equation}
The infimum includes adaptive protocols $\mathcal P$ and arbitrary estimators;
no unbiasedness condition is imposed. Calibration, repetitions, waiting and
failures all share the same deadline. Clipping an estimate to $\Omega$ cannot
increase its squared error. In this model,
\begin{equation}
 R_{\mathrm{fixed}}^*(T;\Omega)=\Theta(T^{-1}),\qquad
 R_{[a,b]}^*(T;\Omega)=\Theta(T^{-2}).
 \label{markov:globaltheorem}
\end{equation}
The interval-rate class and its endpoint-only subclass have distinct
finite-budget risk functions: $R_{[a,b]}^*$ and $R_{\{a,b\}}^*$.
Allowing more rates gives $R_{[a,b]}^*\le R_{\{a,b\}}^*$; equality at a
fixed $T$ is not asserted. Here $\Theta$ means upper and lower bounds by fixed positive multiples for
sufficiently large $T$; the constants carry the required units and depend on
the fixed rates and interval. We next connect this statement to both a
converse for biased estimators and an explicit construction.

The use of a prior below does not change the estimation task into one in which
the frequency is drawn from a known distribution. If an estimator worked
well at every allowed frequency and direction, it would also work well on
average under any weighting of those possibilities. Exhibiting one difficult
weighted average is therefore enough to lower-bound its worst-case error.

Choose an interior half-width $\ell>0$ and center $\omega_c$ with
$[\omega_c-\ell,\omega_c+\ell]\subseteq\Omega$. On a nuisance path use the prior
\begin{equation}
 \pi_\ell(u)=\frac1\ell\cos^2\left(\frac{\pi u}{2\ell}\right),
 \qquad -\ell<u<\ell,\qquad I(\pi_\ell)=\frac{\pi^2}{\ell^2}.
\end{equation}
This is a proof device inside the worst-case parameter family, not information
supplied to the controller. For the joint score
$S=\partial_u\log[p(x\mid u)\pi_\ell(u)]$, integration by parts yields
$\mathbb E[(\widehat\omega-\omega_c-u)S]=1$; the endpoint term vanishes
because the prior density does. The likelihood score has conditional mean
zero, so $\mathbb E S^2=\mathbb E_{\pi_\ell}J_u+\pi^2/\ell^2$.
The Cauchy--Schwarz inequality therefore gives prior-average squared error at
least $1/(\mathbb E_{\pi_\ell}J_u+\pi^2/\ell^2)$, without an unbiasedness assumption.
Worst-case risk is at least this average.

The single-rate cancelling path gives
\begin{equation}
 R_{\mathrm{fixed}}^*(T;\Omega)\ge
 \frac1{T/\gamma+\pi^2/\ell^2}.
 \label{markov:globalfixedlower}
\end{equation}
For every selected rate in $[a,b]$, a supersolution of \eqref{markov:ode} is
$B_T=(rT+\sqrt{CT})^2$: its derivative exceeds
$C+2r\sqrt{B_T}$ by $r\sqrt{CT}\ge0$. Thus
\begin{equation}
 L(T):=\frac1{(rT+\sqrt{CT})^2+\pi^2/\ell^2}
 \le R_{[a,b]}^*(T;\Omega)\le R_{\{a,b\}}^*(T;\Omega).
 \label{markov:globalcontrollower}
\end{equation}
In particular the asymptotic global constants obey
\begin{equation}
 \liminf_{T\to\infty}T R_{\mathrm{fixed}}^*\ge\gamma,\qquad
 \liminf_{T\to\infty}T^2R_{[a,b]}^*\ge\frac1{r^2}.
\end{equation}
These are lower bounds, not claims that the locally optimal coefficients
have globally attainable reciprocals.

For the matching single-rate upper scaling, reference the terminal readout
to the known interval center $\omega_c$ and repeat trials of fixed duration
$\tau=\min\{1/\gamma,\pi/(2W)\}$.
A trial with no error has probability $p_0=e^{-\gamma\tau}$ and produces
phase $\omega\tau$ independently of $\phi$, read relative to the known
phase $\omega_c\tau$. Every trial costs
$\tau$, including those with errors. Among
$m=\lfloor T/\tau\rfloor$ attempts, let $K$ be the number with no error;
$K$ is binomial with parameters $(m,p_0)$. On those useful trials measure the
referenced sine quadrature
$O_c=\cos(\omega_c\tau)Y_L-\sin(\omega_c\tau)X_L$.
Clip the sample mean to the possible sine range and invert
it. Since $W\tau/2\le\pi/4$, the inverse has Lipschitz constant at most
$\sqrt2/\tau$, and conditional risk is at most $2/(K\tau^2)$ for $K>0$.
Use the midpoint if $K=0$. The identities
\begin{equation}
 \mathbb E\frac1{K+1}
 =\frac{1-(1-p_0)^{m+1}}{(m+1)p_0},\qquad
 \frac{\mathbf1_{K>0}}K\le\frac2{K+1}
\end{equation}
then give the uniform bound
\begin{equation}
 R_{\mathrm{fixed}}^*(T;\Omega)
 \le\frac4{\tau^2(m+1)p_0}
     +\frac{W^2}{4}(1-p_0)^m
 =O(T^{-1}).
\end{equation}
This rejection scheme establishes a rate, not an optimal constant. Its
resource count uses attempted trials, so the successful records are not free.

\subsection{Resolve phase ambiguity, then charge the whole random schedule}

There are now two distinct ways an experiment can fail: the measurements
can select the wrong phase turn, or the physical schedule can remain
unfinished at the deadline. The next construction bounds both. Its short
interrogations locate the turn; its long ones refine the position within it.
The physical duration is charged for the complete sequence of attempts.

For two rates, use \eqref{markov:stoppedchannel} to implement exact Ramsey
experiments at effective durations $d_k=2^kd_0$, $k=0,\ldots,K$, where
$Wd_0\le\pi/2$. A known center rotation removes $\omega_c d_k$ from the
readout. At each stage measure both sine and cosine, each on $n_k$ fresh
logical probes. All repetitions are physical experiments. Bounded-outcome
concentration gives
\begin{equation}
 \Pr(\hbox{stage $k$ is inaccurate})\le4e^{-n_k/128}
\end{equation}
when inaccuracy means either quadrature mean differs by more than $1/8$.
Otherwise the inferred phase error is less than $\varepsilon=\pi/6$. Choose
\begin{equation}
 n_k=\left\lceil128\bigl[(K-k)\log8+\log16\bigr]\right\rceil.
 \label{markov:taper}
\end{equation}
The brackets contain natural logarithms. Earlier stages receive more
repetitions because an incorrect coarse fringe is more costly.

Choose at each stage the frequency representative nearest the preceding
estimate. On consecutive accurate stages, a phase error below $\pi/6$ is
small enough to select the correct representative when the duration doubles.
The starting stage is unambiguous on the prescribed band. If the first
inaccurate stage is $k$, the final frequency error is nevertheless
$O(2^{-k}/d_0)$: each subsequent nearest-representative update moves by at
most $\pi/d_\ell$, and those movements have a bounded geometric sum.
Final clipping reduces risk. The stage failure probabilities from
\eqref{markov:taper} are at most $8^{-(K-k)}/4$. Therefore
\begin{equation}
 \sup_{\omega,\phi}\mathbb E(\widehat\omega-\omega)^2
 \le\frac{\varepsilon^2}{d_K^2}
 +\sum_{k=0}^K\frac{9\pi^2}{d_k^2}\,\frac14\,8^{-(K-k)}
 \le\frac{C_{\rm ph}}{d_K^2},
 \label{markov:riskseries}
\end{equation}
where $C_{\rm ph}=\varepsilon^2+9\pi^2/2$ is a sufficient constant.
The first inaccurate stage leaves error at most $3\pi/d_k$; summing
$4^{K-k}8^{-(K-k)}=2^{-(K-k)}$ gives the displayed bound. The total
prescribed effective exposure is
\begin{equation}
 S_K=2\sum_{k=0}^K n_k d_k=\Theta(2^K).
 \label{markov:effectivecost}
\end{equation}
The factor two pays for both quadratures. The taper avoids an extra
logarithmic factor in the total cost: long interrogations require fewer
repetitions than short, coarse ones.

Effective exposure is not yet physical time. Replace each Ramsey operation
in this schedule by a fresh stopped experiment. Its duration law is
independent of phase, direction and readout. Multiplying their transforms
in \eqref{markov:laplace} gives, for the total physical duration,
\begin{equation}
 \mathbb E e^{-q\tau_{\rm total}}=e^{-S_Kk(q)}.
\end{equation}
Choose a fixed $0<\zeta<1$ and the largest $K$ for which
$S_K\le\zeta rT$. The sequence \eqref{markov:effectivecost} grows
geometrically, so this still gives $S_K=\Theta(T)$. For sufficiently small
fixed $\theta>0$, the exponential-moment bound gives
\begin{equation}
 \Pr(\tau_{\rm total}>T)
 \le\exp[-\theta T-S_Kk(-\theta)]\le e^{-cT}
 \label{markov:deadline}
\end{equation}
for some fixed $c>0$ with inverse-time units. Indeed
$-k(-\theta)=\theta/r+O(\theta^2)$, and the strict margin $\zeta<1$
makes the leading exponent negative. Terminate the experiment at $T$ and
use a bounded fallback estimate if incomplete. This contributes at most
$W^2e^{-cT}$ to risk. On completed schedules, the quadrature statistics are
unaffected by conditioning on duration, because completed channels are exact
and duration records are parameter-independent. Equation
\eqref{markov:riskseries} contributes $O(T^{-2})$.

The endpoint construction gives $R_{\{a,b\}}^*(T;\Omega)\le
U_{\rm endpoint}(T)$, where
\begin{equation}
 U_{\rm endpoint}(T)=\frac{C_{\rm ph}}{d_K^2}+W^2e^{-cT}=O(T^{-2})
\end{equation}
for the schedule chosen above and all sufficiently large $T$.
Together with the interval converse,
\begin{equation}
 L(T)\le R_{[a,b]}^*(T;\Omega)\le R_{\{a,b\}}^*(T;\Omega)
 \le U_{\rm endpoint}(T).
\end{equation}
Both selectable classes therefore have order $T^{-2}$, proving
\eqref{markov:globaltheorem}. The argument pays for every repeated probe,
every waiting period and every unfinished schedule. It uses a fixed frequency
band, resolves its phase ambiguity explicitly and does not promote a local
Fisher matrix into a global estimator. Its role beside the finite-collision
theorem is to expose the same resource distinction through an exact
continuous-time mechanism, with its stronger ideal monitoring assumption
kept visible.


\section{What changes when contacts resolve a collision?}
\label{time:access}

The coupling comparison fixes when the controller can act: only at collision
boundaries. That restriction has physical content. A bath qubit becomes
correlated with the sensor during its collision, and a sensor pulse applied
before that collision ends changes their subsequent joint evolution. An
endpoint noise channel does not specify this response to an intervening pulse.
We now keep just one coupling and change temporal access. This is a supporting
comparison, with its own estimation proof; it is not a consequence of the
boundary-only coupling theorem.

Fix $g,h>0$ with $\chi=gh\in(0,\pi/4)$, and retain the same Hamiltonian,
shared direction, protected memory and fixed compensation. The horizon is
$T=nh$. There are at most $n-1$ nonterminal system--memory contacts in total,
of which at most $m$ occur strictly inside cells, where $0\le m\le n-1$.
These are hard bounds on every possible classical history. Boundary contacts
consume the total allowance too. Initial preparation and readout exactly at
$nh$ are free; an earlier readout is a contact and does not shorten the
charged horizon. The permitted controllers may adapt their classical contact
times and instruments, but cannot access the bath or retune the Hamiltonian.

On a sufficiently small fixed band
$\Omega=[\omega_0-E,\omega_0+E]$, the optimal worst-case frequency error obeys
\begin{equation}
 R_{n,m}^{*}(\Omega)
 =\Theta_{\chi,Eh}\!\left(\frac{1}{h^2(n+m^2)}\right).
 \label{time:contactlaw}
\end{equation}
The constants are independent of $n,m$ and the unknown direction. The factor
$h^{-2}$ supplies the units of squared frequency; $\chi$ and $Eh$ are
dimensionless. Thus $m$ of order $\sqrt n$ marks the crossover in order, and
$m$ proportional to $n$ permits inverse-square risk. This does not give an
optimal numerical constant or the largest usable band.

The two terms describe two usable resources. Repeated ordinary one-cell
measurements give the $n$ term. Interior contacts permit a coherent
interrogation whose useful length grows with $m$, giving the $m^2$ term.
The construction below chooses the better of these designs; it need not
combine their information in one experiment. If $m$ grows much more slowly
than $\sqrt n$, the interior contacts do not change the asymptotic order.
If $m$ is proportional to $n$, they do.

\subsection{An exact temporal menu at one strength}

There are two cell settings: leave a cell untouched, or apply one sensor
$Z_S$ pulse at its midpoint. This is an echo: the pulse reverses the noisy
interaction, since $Z_SN_\phi Z_S=-N_\phi$, while leaving the signal term
unchanged, since $Z_SZ_SZ_S=Z_S$. At zero residual frequency the second half
exactly undoes the first half's noisy evolution. Away from zero, we must
calculate the complete evolution rather than assume perfect cancellation.
An additional known boundary operation removes the extra rotation introduced
by the pulse itself; this is the frame correction mentioned below. Let
$R=(g^2+\epsilon^2/4)^{1/2}$ and $r=Rh$. The corrected echoed
propagator, meaning the unitary for the complete cell, is
\begin{equation}
 V=e^{-iH_-h/2}e^{-iH_+h/2},\qquad
 H_\pm=\frac{\epsilon}{2}Z_S\otimes I_E
       \pm gN_\phi\otimes X_E.
\end{equation}
Multiplication gives the parity branches
\begin{align}
 K_{0,e}&=uI-ivZ_S,& K_{1,e}&=-ixN'_\phi,\label{time:echo}\\
 u&=\cos r+\frac{g^2}{R^2}(1-\cos r),&
 v&=\frac{\epsilon}{2R}\sin r,\nonumber\\
 x&=\frac{\epsilon g}{2R^2}(1-\cos r),&
 N'_\phi&=\partial_\phi N_\phi .\nonumber
\end{align}
Here $u,v,x$ are dimensionless real numbers and $u^2+v^2+x^2=1$.
On an odd result, an additional known $Z_S$ within that same boundary
instrument uses $Z_SN'_\phi=-iN_\phi$ to make the flip operator proportional
to $N_\phi$. The frame correction, parity measurement, conditional correction
and next preparation are components of one instrument at one time; the
midpoint pulse is a separate contact.

Both settings consequently have the same unknown normalized flip, while
their no-flip phases differ. Denote the untouched setting by $p$ and the
echoed setting by $e$. Their transition probabilities and phase derivatives
at compensation are
\begin{equation}
 p_e(0)=0,\quad p_p(0)=\sin^2\chi,\qquad
 \theta'_e(0)=A:=\frac{\sin\chi}{g},\quad
 \theta'_p(0)=d_p:=\frac{\tan\chi}{g}.
 \label{time:echoderivatives}
\end{equation}
Primes on scalar cell functions mean differentiation with respect to the
residual frequency $\epsilon$; $N'_\phi$ retains its explicit meaning
$\partial_\phi N_\phi$. Thus $A$ and $d_p$ have units of time. The zero
value of $p_e$ holds at compensation only. Away from it, $p_e=x^2$ is usually
positive.

Use an echoed cell after an even parity result and a plain cell after an odd
result. Every record ending even again has exact cancellation of $\phi$ and
a known signed sum of the no-flip phases. To verify usable sensitivity,
freeze the two phase rewards at $A$ and $d_p$. Their stationary signed mean is
\begin{equation}
 \mu^0(\epsilon)=
 \frac{p_p(1-p_e)A-p_e(1-p_p)d_p}{p_e+p_p},
 \qquad \mu^0(0)=A>0.
\end{equation}
Continuity supplies one fixed band on which $\mu^0\ge3A/4$,
$p_e<p_p<1/2$, and each actual derivative differs from its frozen value by at
most $A/4$. Also $p_e+p_p$ stays bounded below, uniformly controlling how
rapidly dependence between well-separated parity outcomes decays.
The record gate already developed for the coupling menu then accepts at least
one quarter of sufficiently long blocks and gives, on every accepted record
of length $L$,
\begin{equation}
 AL/4\le \Phi'_{\mathcal R}(v)\le \beta L
 \quad (|v|\le E)
\end{equation}
for a fixed finite time constant $\beta$. A positive lower bound on $p_e$
is unnecessary: at compensation the process never leaves its initial even
sector. Plain
one-cell trials provide the paid initial localization. The same conditional
inversion and multiscale schedule therefore give risk $C_1/(h^2n_{\mathrm{use}}^2)$ with
at most $n_{\mathrm{use}}$ used cells, including calibration, repetitions and rejected
records.

For $n\ge3$ and $m\ge1$, choose
\begin{equation}
 n_{\mathrm{use}}=\min\!\left\{m,\left\lfloor\frac{n-1}{2}\right\rfloor\right\}.
 \label{time:allocation}
\end{equation}
The whole estimator uses at most $n_{\mathrm{use}}$ midpoint pulses and $n_{\mathrm{use}}$ boundary
instruments. Its last readout is counted because it occurs before $nh$.
Thus $2n_{\mathrm{use}}\le n-1$ and $n_{\mathrm{use}}\le m$. Store the resulting classical estimate
through the remaining charged cells. For example, $n=101,m=20$ permits
$n_{\mathrm{use}}=20$, uses at most 40 contacts, and still pays for all 101 cells. It does
not receive another 101 cells for calibration.

Since $n_{\mathrm{use}}\ge m/3$, this construction has risk at most
$9C_1/(h^2m^2)$. A separate design uses all $n$ cells for plain-cell
calibration and has risk at most $C_0/(h^2n)$ with $n-1$ boundary contacts.
Choose the better design from the known budgets. The elementary inequality
\begin{equation}
 \min\!\left\{\frac{C_0}{n},\frac{9C_1}{m^2}\right\}
 \le \frac{C_0+9C_1}{n+m^2}
\end{equation}
proves the upper order in \eqref{time:contactlaw}; the plain design covers
$m=0$ and the small values of $n$.

\subsection{Why adaptivity does not evade the contact bound}

The proof uses the geometry of small changes in a quantum state. Represent
the response to a small parameter change by a vector. Contributions from
untouched cells are perpendicular in this mathematical state space, so their
squared lengths add, as in Pythagoras' theorem. Contributions from contacted
cells may instead align, so their lengths can add. This is the origin of the
two scales $\sqrt n$ and $m$ in the information bound.

Use the plain-cell nuisance curve
$\phi'=q_p\theta'_p/(2p_p)$ from the fixed-coupling converse. It is a
mathematical subfamily used to lower-bound worst-case error, not information
supplied to the controller. Put
\begin{equation}
 p_{\min}=\chi^2\left(\frac{\sin r_*}{r_*}\right)^2,\quad
 K_\phi=\sup_{|\epsilon|\le E}|\phi'(\epsilon)|,\quad
 a_0=\frac{h}{2\sqrt{p_{\min}}},\quad
 b_0=h\sqrt{\frac14+g^2K_\phi^2},
\end{equation}
where $r_*=(\chi^2+(Eh/2)^2)^{1/2}<\pi/4$ bounds $Rh$ on the band.
Both $a_0,b_0$ have units of time. A cell with no interior contact has
zero derivative overlap and derivative norm at most $a_0$ in the chosen
purification. A cell with interior contacts instead has derivative norm
at most $b_0$, because along the same curve
\begin{equation}
 \dot H=\tfrac12Z_S\otimes I_E
       +g\phi' N'_\phi\otimes X_E,
 \qquad \|\dot H\|=\sqrt{\tfrac14+g^2(\phi')^2}.
\end{equation}
The norm is the operator norm, the largest vector-stretching factor. The equality
uses anticommutation. Integrating the unitary derivative over the cell's
physical duration $h$ gives the bound even with intermediate adaptive
instruments.

The hard history constraint is essential to the rest of the argument.
Dilate the instruments---represent them as unitary operations on a larger
state space---keeping a mathematical archive of which cells
contained an interior contact. This archive is not a new observation
available to the controller. Derivative contributions from untouched cells
are mutually orthogonal, so their total norm is at most $a_0\sqrt n$.
For contacted cells, let $D_j$ be the vector contribution from cell $j$.
The projector $Q_z$ selects the part associated with one archived activity
history $z$, where $z_j=1$ marks a contacted cell. Every history
contains at most $m$ active cells. Cauchy--Schwarz within each history gives
\begin{equation}
 \left\|\sum_jD_j\right\|^2
 =\sum_z\left\|\sum_{j:z_j=1}Q_zD_j\right\|^2
 \le m\sum_j\|D_j\|^2\le m^2b_0^2.
\end{equation}
The last step charges the expected number of active cells by the hard bound
$m$. Merely summing square roots of separate activity probabilities would
lose this dependence. The entrance activity flags are legitimate because
operations on protected memory alone commute with the cell evolution;
they can be moved before it until the first actual system contact.

Consequently the complete accessible information along the curve obeys
\begin{equation}
 J_\epsilon\le4(a_0\sqrt n+b_0m)^2
 \le4(a_0^2+b_0^2)(n+m^2).
\end{equation}
A fixed cosine-squared prior of half-width $E$ yields the global bound
\begin{equation}
 R_{n,m}^{*}(\Omega)\ge
 \frac{1}{4(a_0^2+b_0^2)(n+m^2)+\pi^2/E^2}.
\end{equation}
This applies to biased estimators and all retained branches. Clipping and
finite quantization transfer the bounded-loss statement to the stated
operational limits; continuity of Fisher information is not assumed.
The converse may grant extra boundary operations as a relaxation, but never
relaxes the hard interior budget. Together with the construction it proves
\eqref{time:contactlaw}.

\section{Timing knowledge changes the experiment}
\label{time:timing}

A pulse being displaced by a known amount and a pulse having an unobserved
random displacement are different statistical experiments, even at the same
error magnitude. A simple timing example makes the distinction concrete.
Divide a cell into 100 equal clock ticks, so that its midpoint is tick 50:

\begin{center}
\begin{tabular}{@{}p{0.23\linewidth}p{0.70\linewidth}@{}}
\toprule
Timing information & What the controller knows\\
\midrule
Known fixed offset & Every pulse arrives at tick 55, and this fact is supplied.\\
Unknown fixed offset & Every pulse has the same displacement, but its value must be learned. This calibration problem is not solved here.\\
Hidden random jitter & Each pulse independently arrives at tick 45 or 55; the probabilities are known, but the realized choice is never reported.\\
\bottomrule
\end{tabular}
\end{center}
The ticks are only an illustration, not an additional measurement available
to the controller. To isolate the timing question, keep one coupling and permit
only plain cells or cells with one $Z_S$ pulse. A displaced pulse occurs at
$h/2+\delta_t$, with $|\delta_t|<h/2$; there is no further interior contact in
that cell. Here $\delta_t$ is a timing displacement, distinct from the
Markovian half-rate contrast $\delta=(b-a)/2$. Its corrected propagator is
\begin{equation}
 V_{\delta_t}=e^{-iH_-(h/2-\delta_t)}e^{-iH_+(h/2+\delta_t)}.
\end{equation}
Its known pulse frame is included in the counted endpoint instrument.
The exact branch coefficients are
\begin{align}
 D&=\cos(2R\delta_t)-\cos(Rh),&
 u&=\cos(Rh)+\frac{g^2}{R^2}D,\nonumber\\
 v&=\frac{\epsilon}{2R}\sin(Rh),&
 w&=\frac{g}{R}\sin(2R\delta_t),\nonumber\\
 x&=\frac{\epsilon g}{2R^2}D,\label{time:displacedcoeffs}
\end{align}
with
\begin{equation}
 K_0=uI-ivZ_S,\qquad K_1=-i(wN_\phi+xN'_\phi),
 \qquad u^2+v^2+w^2+x^2=1.
 \label{time:displacedbranches}
\end{equation}
Each branch is proportional to a unitary because $N_\phi$ and $N'_\phi$
anticommute. The odd branch has a frequency-dependent axis, but for a supplied
known $\delta_t$ that dependence is known.

\subsection{A supplied fixed displacement preserves a usable phase}

Fix $0<|\delta_t|<h/2$ throughout the experiment and suppose its value is part
of the device specification. Near compensation, $w$ has fixed nonzero sign.
Absorb that sign into an irrelevant scalar phase and define
\begin{equation}
 \theta_e=2\arctan(v/u),\qquad
 \zeta_e=\arctan(x/w).
\end{equation}
The normalized odd branch is then $N_{\phi+\zeta_e}$ up to a scalar phase.
For a plain cell, $\zeta_p=0$. Using pulsed cells in the even sector and plain
cells in the odd sector gives the exact final-even phase
\begin{equation}
 \Phi_{\mathcal R}(\epsilon)=
 \sum_i\sigma_{i-1}\bigl[
  \mathbf1_{\rm stay}\theta_{j_i}(\epsilon)
  -2\mathbf1_{\rm flip}\zeta_{j_i}(\epsilon)\bigr].
 \label{time:knownphase}
\end{equation}
The extra flip-axis terms must be retained. They are known functions of
frequency, while the common unknown angle still cancels in pairs.

For a direct sensitivity check, set $z=2g\delta_t$, so $0<|z|<\chi$.
At zero residual frequency,
\begin{align}
 p_e&=\sin^2z,&p_p&=\sin^2\chi,\nonumber\\
 d_e&=\theta'_e(0)=\frac{\sin\chi}{g\cos z},&
 z_e&=\zeta'_e(0)=\frac{\cos z-\cos\chi}{2g\sin z}.
\end{align}
Put $f_e=(1-p_e)d_e-2p_ez_e$ and
$f_p=(1-p_p)\tan\chi/g$. The exact contrast is
\begin{equation}
 \frac{f_e}{p_e}-\frac{f_p}{p_p}
 =\frac{(\sin\chi-\sin z)\sin(\chi+z)}
 {g\sin\chi\,\sin^2z}>0.
\end{equation}
All factors determining the sign are positive for $-\chi<z<\chi$ except
at the excluded zero. Thus the stationary derivative drift
$(p_pf_e-p_ef_p)/(p_e+p_p)$ is strictly positive; denote it by $\mu_{\delta_t}$.
The record score must now include the new flip contributions:
\begin{equation}
 D_L^0=\sum_{i=1}^L\sigma_{i-1}
 \left[\mathbf1_{\rm stay}\theta'_{j_i}(0)
       -2\mathbf1_{\rm flip}\zeta'_{j_i}(0)\right].
\end{equation}
It is computed from the observed record and the supplied displacement.
Its stationary mean is $\mu_{\delta_t}$. Accept even-ending records with
$D_L^0\ge\mu_{\delta_t} L/2$. All these one-cell rewards are bounded for the
fixed nonzero displacement. The same two-state mean and variance argument
therefore gives a positive acceptance probability for sufficiently long
blocks. Continuity on a sufficiently small fixed band keeps each derivative
close to its frozen value, certifying a whole-band positive phase slope.
The same conditional estimator then gives inverse-square risk.
The band and constants may depend on the supplied displacement.
Using \eqref{time:allocation} charges at most $n_{\mathrm{use}}$ interior and $n_{\mathrm{use}}$ boundary
contacts, so this yields the quadratic global order when $m=\Theta(n)$.
The ideal $\delta_t=0$ case is already covered by the midpoint construction.

This argument does not need to retime the pulse. If calibrated retiming is
allowed, shifting the command by $-\delta_t$ restores the midpoint instead.
Neither statement supplies a method for learning an unknown $\delta_t$.
Such learning would need its own identifiable likelihood, accuracy analysis
and charged calibration cost.

\subsection{Hidden independent jitter mixes the odd branch}

Now let each pulsed cell independently have $\delta_t=+d$ or $-d$ with equal
probability, where the amplitude $0<d<h/2$ is known but its realized signs are
never reported. The signs are independent of the unknown parameters, prior
records and controls. At cell entrance the controller can choose plain or
jittered, but has no further pulse-position menu, extra interior pulse,
coupling choice or physical compensation change.

Only $w$ in \eqref{time:displacedcoeffs} changes sign. Averaging the two
physical channels gives three Kraus operators, evaluated at $\delta_t=d$:
\begin{equation}
 L_0=uI-ivZ_S,\qquad L_1=-iwN_\phi,\qquad L_2=-ixN'_\phi.
 \label{time:jitterkraus}
\end{equation}
Parity distinguishes $L_0$ from the combined odd branch. It does not reveal
which odd operator acted. A putative record identifying $L_1$ or $L_2$ would
supply information absent from this experiment.

For input state $\rho$, the normalized state after an odd result is
\begin{equation}
 \rho_{\rm odd}=
 \frac{w^2N_\phi\rho N_\phi+x^2N'_\phi\rho N'_\phi}{w^2+x^2}.
\end{equation}
Two alternatives remain unresolved. In general this is a mixture, not one
known unitary rotation whose phase can be calculated from the record.
That explains why the earlier estimator does not directly apply. A separate
converse is still needed to rule out other permitted strategies.

For every fixed $d>0$, this restricted menu has
\begin{equation}
 R_{n,m}^{*}(d;\Omega)=\Theta(n^{-1}),
 \qquad 0\le m\le n-1,
 \label{time:jitterrisk}
\end{equation}
on a fixed identifiable band, uniformly over the shared direction. Constants
may depend on $g,h,E,d$. The upper bound simply uses plain-cell calibration;
it fits every $m$ and costs all $nh$. The limitation requires a new converse,
because good local sensitivity at compensation still survives.

Choose a fixed compact residual-frequency interval strictly inside $(0,E)$.
Use the plain-cell nuisance curve $\phi'=k(\epsilon)$ from the earlier
converse. For the jittered cell write $p=w^2+x^2$ and $a=uv'-vu'$. Along this
curve,
\begin{equation}
 i\sum_jL_j^\dagger\dot L_j=(a-pk)Z_S,
 \qquad L_1^\dagger L_2=iwxZ_S.
\end{equation}
On the chosen positive interval $w,x>0$. The same accessible channel has
several equivalent descriptions using inaccessible environmental labels.
Changing that description leaves every accessible measurement probability
unchanged. We choose one convenient description to prove a bound; this choice
is called an environment gauge and grants no control over the physical bath.
Its smooth unitary change of basis has a Hermitian derivative generator $G$,
a $3\times3$ matrix acting on the three Kraus labels, with entries
\begin{equation}
 G_{12}=i\eta_G,\qquad G_{21}=-i\eta_G,\qquad
 \eta_G=\frac{a-pk}{2wx},
\end{equation}
and all other entries zero. The entries of $G$ have units of time because
the derivative is with respect to frequency. In the derivative convention
$\widetilde{\dot L}=\dot L-iGL$, $L$ denotes the column of three Kraus
operators. This choice cancels the overlap exactly:
\begin{align}
 i\sum_jL_j^\dagger\widetilde{\dot L}_j&=0,\nonumber\\
 \sum_j\widetilde{\dot L}_j^\dagger\widetilde{\dot L}_j
 &=\bigl[(u')^2+(v')^2+(w')^2+(x')^2\nonumber\\
 &\hspace{2.1em}+(wk+\eta_Gx)^2+(xk+\eta_Gw)^2\bigr]I.
 \label{time:gaugealpha}
\end{align}
For completeness the unitary basis change exists by solving
$U'_E=-iU_EG$ on that interval. It is an analytical freedom in a purification,
not bath access granted to the experimenter. The scalar in
\eqref{time:gaugealpha}, and the corresponding plain-cell scalar, have one
finite bound $A_*$ on this compact interval.

Both selectable cells now have zero overlap on the same nuisance curve.
Successive derivative contributions are therefore orthogonal even with
adaptive setting choices, arbitrary protected memory and all branches
retained. The complete information is at most $4nA_*$. A cosine-squared
prior of fixed half-width $\ell$ contained in the positive interval gives
\begin{equation}
 R_{n,m}^{*}(d;\Omega)\ge
 \frac{1}{4nA_*+\pi^2/\ell^2}.
\end{equation}
The prior is deliberately away from compensation, where $x=0$ and the
chosen gauge would be singular. A worst-case risk over the entire band is
at least this fixed-prior average. Thus the singularity at zero does not
invalidate the global limitation, and the proof makes no pointwise risk
claim at every frequency. The bounded-loss limit argument again extends the
result to the allowed operational limits.

\subsection{Local information and a necessary accuracy scale}

At compensation the third Kraus operator in \eqref{time:jitterkraus}
vanishes. Near zero it has the form $L_2(\epsilon)=\epsilon B+O(\epsilon^2)$
for a fixed operator $B$ with units of time. Its contribution to the state is consequently
$\epsilon^2B\rho B^\dagger+O(\epsilon^3)$: it does not affect the first
derivative exactly at zero, although it matters away from zero.
The first-order instrument therefore still
permits a positive phase drift. Put $z_d=2gd$, $A=\sin\chi/g$,
$p_e=\sin^2z_d$ and $p_p=\sin^2\chi$. The even-jittered/odd-plain protocol
has local signed-phase drift per physical time
\begin{equation}
 v_d=\frac{A[\sin^2\chi\cos z_d-\sin^2z_d\cos\chi]}
 {h(p_p+p_e)}>0.
\end{equation}
Positivity follows from the strict increase of $\sin^2z/\cos z$ on
$(0,\pi/4)$ and $0<z_d<\chi$.
Here is the finishing operation used for the local claim. For a sufficiently
large allocation of $N$ cells, put $q_p=1-\sin^2\chi<1$ and reserve
$k_N=\lceil3\log N/(-\log q_p)\rceil$ of those cells. Run the parity policy
for the first $N-k_N$ cells. If parity is even, read out immediately; if it
is odd, use plain cells until it first becomes even or the allocation ends.
At an even endpoint measure $Y_L$, retain the record and store the result.
An unfinished trial retains a failure flag. The return failure probability
at compensation is at most $q_p^{k_N}\le N^{-3}$. The reserved number is only
$O(\log N)$, so it does not change the leading accumulated phase derivative.
The complete record with this readout has pointwise classical score information
$v_d^2(Nh)^2+o(N^2)$ for $N$ used cells. Allocating
$N=\min\{m,\lfloor(n-1)/2\rfloor\}$ pays the early readout and fits the
original contact budget. If $m=\Theta(n)$, quadratic local information and
the global inverse-linear law \eqref{time:jitterrisk} therefore coexist.
A particularly sensitive readout at one special frequency need not work
uniformly well across an interval. Here the local statement concerns the
classical outcome probabilities of a specified readout, not an assumption
that quantum Fisher information remains continuous when a noise branch
vanishes.

The same away-from-zero gauge supplies a necessary condition when jitter
shrinks with the budget. Fix $d_{\max}<h/2$ and put $\xi=d/h$. Uniformly for
$0<d\le d_{\max}$ on the fixed prior interval, $w$ is bounded below by a
constant times $\xi$, while $x$ stays bounded below by a positive constant.
The gauged derivative cost is consequently at most $C h^2/\xi^2$ per
pulsed cell and $C h^2$ per plain cell. Charging at most $m$ pulsed cells
gives, with a fixed constant $C$,
\begin{equation}
 R_{n,m}^{*}(d;\Omega)\ge
 \max\!\left\{
 \frac{1}{n^2h^2+\pi^2/\ell^2},\,
 \frac{1}{Ch^2(n+m/\xi^2)+\pi^2/\ell^2}
 \right\}.
\end{equation}
The first term is the independent signal-generator limitation. For
$m=\Theta(n)$, the second denominator can grow quadratically only if
$1/\xi^2$ grows at least linearly. Risk of order $n^{-2}$ would therefore require
\begin{equation}
 d=O(h/\sqrt n).
\end{equation}
The reference scale $h/\sqrt n$ halves when $n$ quadruples. This illustrates
the necessary asymptotic order, not an exact finite-budget tolerance.
Meeting that scale alone does not
establish a working estimator: necessity rules out inadequate accuracy,
whereas sufficiency would require a construction that succeeds.
This is not an optimized joint law
in $n,m,d$, and does not cover recorded timing errors, arbitrary pulse
sequences, finite-duration controls or additional coupling settings.

\section{Which idealizations carry the conclusion?}
\label{time:boundaries}

The shared transverse direction is essential. Two applications of the
same error cancel, and the same unknown angle constrains the nuisance
curve in both settings. If a continuous-time setting $j$ has an independent
unknown angle $\phi_j$, signal coefficient $s_j$ and positive rate
$\gamma_j$, the path $\phi'_j=s_j/(2\gamma_j)$ cancels the coherent
contribution separately in every setting. It gives the all-strategy ceiling
\begin{equation}
 \mathcal I_\omega^{*}(T)\le
 T\max_j\frac{s_j^2}{\gamma_j}.
\end{equation}
Known angular offsets between settings are not independently unknown axes.
Likewise a rate knob must have its stated signal response. If setting $j$
adds an unconstrained unknown frequency offset $\delta_j$, then
$(s_j\omega+\delta_j)Z/2$ is invariant under
$\omega\mapsto\omega+u$ and $\delta_j\mapsto\delta_j-s_ju$.
The target is then unidentifiable without an additional constraint.

An uncorrectable parallel noise component gives a second sharp boundary.
In the continuous model add the independent unobserved operator
$L_z=\sqrt\kappa Z$, with known $\kappa>0$. Its contribution is
$\kappa(Z\rho Z-\rho)$, so off-diagonal coherence decays as
$e^{-2\kappa T}$. Parity does not identify these errors. The signal lies in
the same operator direction as this noise, and a channel gauge coupling the
identity to $L_z$ with coefficient $-1/(4\sqrt\kappa)$ gives
\begin{equation}
 \mathcal I_\omega^{*}(T)\le\frac{T}{4\kappa}.
\end{equation}
This familiar channel-bound mechanism applies even if the transverse axis
is known. Hence the strict quadratic asymptotic claim is not protection
against arbitrary added noise.

A finite advantage can nevertheless remain. For the prescribed fixed-duration
random-rate readout with a recorded uniform equatorial measurement angle,
parallel dephasing multiplies every entry of its frequency--direction Fisher
matrix by
\begin{equation}
 w(T)=1-\sqrt{1-e^{-4\kappa T}}.
\end{equation}
The nuisance-adjusted information is multiplied by the same factor. For
$b=3a$, lower-rate probability $3/4$, and $T=10/a$, the ideal information is
$200/(7a^2)$. At $\kappa/a=0.01$ it becomes approximately $12.16635/a^2$,
still above the universal ideal single-rate-$a$ bound $T/a=10/a^2$.
This is a finite local-information comparison for a specified readout, not
a global risk ratio or a general hardware error tolerance. The random
physical duration of a stopped interrogation would require its own
averaging; this fixed-duration formula cannot be transplanted to it.

Finite-rate monitoring also needs its own scope. In the continuous Markov
model, parity checks separated by a fixed interval $\Delta$ miss pairs of
jumps. Writing $p_\gamma=(1-e^{-2\gamma\Delta})/2$, the even/odd rate policy
has mean signed signal rate
\begin{equation}
 v_\Delta=\frac{p_ap_b}{\Delta(p_a+p_b)}
 \left(\frac1a-\frac1b\right)>0.
\end{equation}
At compensation, a retained-even logical $Y_L$ readout gives local information
$v_\Delta^2T^2+o(T^2)$. Use the same kind of logarithmic finishing window
within the allocated intervals, retaining the high rate in the odd sector
until an even boundary result, and keep deadline failures as flags.
Thus both the finishing time and failures are charged. Away from
compensation, unobserved residence times form a mixture of phases. This
local finite-cadence result does not prove the continuous-monitoring global
risk theorem at finite cadence. The microscopic global theorem avoids this
inference by analyzing its own exact recorded instrument.

Finally, the charged resource is exposure time under specified ideal control.
Memory is protected, instruments and pulses are instantaneous, the bath cells
are prepared as stated, and the shared axis is stable throughout the run.
Control work, bath preparation expense, pulse duration, readout errors and
unknown timing-offset calibration are not assigned free solutions. Exact
conditional identities and analytic bounds establish the mathematical
results in this resource model; finite numerical checks verify selected
identities and examples. Neither those checks nor a completed exposition
constitutes an independent experimental or expert certification.


\section{What the comparison establishes}
\label{end:synthesis}

The question was how accurately to measure an initially unknown frequency
when the same unknown transverse error keeps acting on the sensor. The
quantity being optimized is mean-square error, including every experimental
outcome, at the hardest frequency and direction in the promised set. The
resource is the full physical duration $nh$, with the stated restrictions
on when the controller can act. The frequency interval, collision duration,
couplings and physical compensation are fixed as $n$ grows.

With one positive coupling, a change in frequency can be partly hidden by
a change in the unknown error direction. The proof chooses a whole curve
of such difficult parameter pairs and shows that no allowed boundary-control
strategy can distinguish them faster than a linear information law permits.
A guarantee at every parameter pair must also hold on average over that
curve, which turns the information limit into an inverse-linear lower bound
on error. Repeated, fully charged one-cell calibration reaches that order.
This is a limitation of all allowed strategies, not just a failure of one
chosen recovery procedure.

With two distinct positive couplings, the construction uses the weaker one
in the even sector and the stronger one in the odd sector. The unknown
direction cancels on every even-ending record. Some such records have little
useful phase left, so a computable rule accepts only those with a certified
response throughout the frequency band. The retained path supplies a known
prediction curve for every candidate frequency; it need not reveal the true
frequency on its own.

The estimator then pays for a coarse estimate and uses progressively longer
blocks to narrow the candidate interval. It compares measured outcomes with
the curves from the paths actually observed. Early mistakes receive greater
protection because they would cause larger errors; all rejected blocks still
consume their full time. The resulting worst-case mean-square error has order
$1/n^2$. A separate bound on the total signal response rules out a better
order. Thus both sides of the main comparison have a matching construction
and converse. Updating the estimate never retunes the physical compensation.

The continuous-time model reaches the same distinction under its own,
stronger assumption of ideal continuous error monitoring. Its first
explanation is statistical: two noise rates expose different combinations of
the same unknown frequency and direction. Its second explanation pairs errors
and stops when the signed signal clock reaches a chosen value. This makes an
exact sensing operation, but the physical waiting time must still be charged.
The finite-cell theorem does not depend on approximating its hardware by
that continuous model. Sharp local information coefficients are not asserted
to be optimal constants in the global estimation error.

Access during a collision is another resource. At one fixed coupling,
the companion contact law combines an ordinary $n$ scale with a coherent
$m^2$ scale, where $m$ bounds the allowed interior contacts. Both its
construction and its limitation use the full contact budget. The same
recorded-phase estimator is useful there, but the physical operation enabling
it is different.

Finally, timing knowledge changes what the record tells us. A supplied
fixed displacement can be included when calculating each path's phase.
Hidden independent jitter can instead leave two unresolved operations after
the same observed parity result. For the stated plain/one-pulse menu,
a global lower bound then restores inverse-linear error at any fixed
nonzero jitter amplitude. High sensitivity at the special compensation point
does not contradict poor worst-case precision over a whole interval.
The shrinking-jitter accuracy condition is necessary only: meeting it has
not been shown sufficient.

The stable shared axis, protected memory, ideal permitted instruments and
fixed band remain essential. Within that model, controlling noise strength
changes the attainable order of global estimation error, with every trial
and failure charged.

\end{document}
